% ==================================================================
%  TESI MAGISTRALE - Mauro Segreti
%  Versione con il RESTYLE di Luca + tutte le CORREZIONI del PDF
%  ------------------------------------------------------------------
%  I box rossi "DA FARE" segnalano le note di Luca che richiedono
%  lavoro tuo (nuove analisi, nuovi capitoli, ecc.).
%  Per farli sparire prima della consegna: \showtodostrue -> false
% ==================================================================
\documentclass[11pt,a4paper,twoside,openright]{report}
\PassOptionsToPackage{export}{adjustbox}

% --- Codifica, font, lingua ---
\usepackage[utf8]{inputenc}
\usepackage{lmodern}   % font vettoriale: serve a microtype con T1
\usepackage[T1]{fontenc}
\usepackage[english]{babel}
\usepackage{microtype}

% --- Matematica ---
\usepackage{amsmath, amssymb, amsfonts, amsthm}
\usepackage{bm}

% --- Figure e tabelle ---
\usepackage{graphicx}
\usepackage{adjustbox}
\usepackage{float}
\usepackage{booktabs}
\usepackage{caption}
\usepackage{subcaption}
\usepackage{lettrine}
\usepackage{tabularx}
\usepackage{array}
\graphicspath{{figures/}{TESI/}{TESI/Figures/}}

% --- DIDASCALIE IN ROSSO ---------------------------------------------
% Per tornare al nero prima della consegna: commenta le due righe qui sotto.
%\captionsetup{font={small,color=red}, labelfont={small,bf,color=red}}
%\captionsetup[sub]{font={small,color=red}, labelfont={small,bf,color=red}}

% --- Layout di pagina ---
\usepackage{geometry}
\geometry{a4paper, top=3cm, bottom=3cm, left=3.5cm, right=2.5cm} % margini per rilegatura
\usepackage{emptypage}
\usepackage[pagestyles]{titlesec}   % NON caricare fancyhdr: sono incompatibili

% --- Colori ---
\usepackage{xcolor}
\definecolor{sapphire}{rgb}{0.03, 0.15, 0.4}

% --- Box colorati ---
\usepackage[most]{tcolorbox}

% --- Link, bibliografia, riferimenti incrociati ---
% Ordine obbligatorio: hyperref -> biblatex -> cleveref
\usepackage{hyperref}
\hypersetup{
    colorlinks=true,
    linkcolor=sapphire,
    citecolor=sapphire,
    urlcolor=sapphire
}
\usepackage[numbers,sort&compress]{natbib}
\usepackage{cleveref}

% --- Formattazione titoli dei capitoli ---
\titleformat{\chapter}[display]
  {\normalfont\bfseries\color{sapphire}}
  {\filleft\huge\chaptertitlename\ \thechapter}
  {20pt}
  {\titlerule\vspace{1ex}\raggedright\Huge}
  [\vspace{1ex}\titlerule]

% --- Intestazioni di pagina ---
\newpagestyle{main}[\normalsize\scshape]{
  \sethead[\thepage][][\chaptertitlename\ \thechapter.\ \chaptertitle]
          {\toptitlemarks\thesection\quad\sectiontitle}{}{\thepage}
  \setheadrule{0.4pt}
}
\pagestyle{main}

% ==================================================================
%  BOX PROMEMORIA (note di Luca da lavorare)
%  Metti \showtodosfalse per nasconderli tutti nella versione finale
% ==================================================================
\newif\ifshowtodos
\showtodofalse
\newcommand{\faExclamation}{}
\newtcolorbox{todobox}[1][Nota di Luca]{
  enhanced, breakable,
  colback=red!4!white,
  colframe=red!65!black,
  coltitle=white,
  fonttitle=\bfseries\sffamily\small,
  fontupper=\small\sffamily,
  title={\faExclamation DA FARE --- #1},
  left=3mm, right=3mm, top=2mm, bottom=2mm,
  boxrule=0.8pt, arc=2mm
}
\newenvironment{todo}[1][Nota di Luca]
  {\ifshowtodos\begin{todobox}[#1]\else\setbox0\vbox\bgroup\fi}
  {\ifshowtodos\end{todobox}\else\egroup\fi}

% ==================================================================
%  MACRO PERSONALI (sempre dopo tutti i \usepackage)
% ==================================================================
\newcommand{\ccterm}[1]{%
  \overline{q}\gamma^\mu\left\{ig_w\frac{#1}{\sqrt{2}}\,\frac{1-\gamma^5}{2}\right\}q}

\begin{document}
% ==================================================================
%  FRONTESPIZIO
% ==================================================================
\begin{titlepage}
    \centering
    \vspace*{1cm}

    % ---------- Tor Vergata ----------
    \includegraphics[width=\textwidth]{logo_tor_vergata.png}

    \vspace*{\stretch{1.5}}

    % ---------- Dipartimento ----------
    {\large\color{sapphire}\textls[180]{\textsc{Department of Physics}}\par}
    \vspace{0.45cm}
    {\large Master's Degree in Physics\par}
    \vspace{0.15cm}
    {\small\itshape Curriculum in Fundamental Interactions and Experimental Techniques\par}

    \vspace*{\stretch{1.7}}

    % ---------- Titolo ----------
    {\color{sapphire}\rule{\textwidth}{0.8pt}}\\[1.1cm]
    {\huge\bfseries\color{sapphire}
     Performance of the ATLAS Muon \\[0.25cm]
     Spectrometer in Run 3 of the LHC\par}
    \vspace{1.1cm}
    {\color{sapphire}\rule{\textwidth}{0.8pt}}

    \vspace*{\stretch{1.7}}


    % ---------- Relatori e candidato ----------
    \begin{minipage}[t]{0.48\textwidth}
        \begin{flushleft}\large
            \emph{Supervisors:}\\[0.15cm]
            Prof. Anna \textsc{Di Ciaccio}\\
            Dr. Marco \textsc{Vanadia}\\
            Dr. Luca \textsc{Martinelli}
        \end{flushleft}
    \end{minipage}
    \hfill
    \begin{minipage}[t]{0.44\textwidth}
        \begin{flushright}\large
            \emph{Candidate:}\\[0.15cm]
            Mauro \textsc{Segreti}
        \end{flushright}
    \end{minipage}

    \vspace*{\stretch{1}}

    % ---------- Data ----------
    {\large Academic Year 2025/2026\par}
    \vspace{0.2cm}
    {\small\color{sapphire}Rome, Italy\par}
\end{titlepage}
% ==================================================================
%  ABSTRACT
% ==================================================================
\begin{abstract}
This thesis presents muon momentum resolution and performance studies for the ATLAS detector, based on the \textit{pp} collision data sample produced at \(\sqrt{s} = 13.6\ \text{TeV}\) at the LHC during Run~3, corresponding to an integrated luminosity of \(332\ \text{fb}^{-1}\) \cite{LuminosityEra,ATLAS:LumiRun3}. The muon momentum resolution are measured using simulated samples of $Z\rightarrow\mu^+\mu^-$ and $Z'\rightarrow\mu^+\mu-$ events. "DA RIVEDERE"
\end{abstract}

%CORRETTO--------------------------------------------------
% CORREZIONE (p.3): 13 TeV -> 13.6 TeV, energia del Run 3.
% CORREZIONE (p.3): 139 fb^-1 e' la luminosita' del Run 2 (140 dopo il ricalcolo).
%   Nel Run 3 ATLAS ha registrato 332 fb^-1. Vedi:
%   https://atlas.web.cern.ch/Atlas/GROUPS/DATAPREPARATION/PublicPlots/2026/DataSummary/figs/intlumivstimeRun123.png
% CORREZIONE (p.3): "This paper" -> "This thesis" ("ah si? e' diventato un paper questa tesi? ahahah")

% ==================================================================
%  INDICE
% ==================================================================
\cleardoublepage
\tableofcontents


% ==================================================================
%  CAPITOLO 1 - INTRODUZIONE
% ==================================================================
\chapter{Introduction}

\lettrine[lines=3,loversize=0.1,lraise=0.05]{T}{} he third data-taking period of the Large Hadron Collider, known as Run 3, started on 5 June 2022, at the end of two-year period of Long Shutdown 2, during which both the accelerator and the experiments underwent an extensive upgrade programme. Protons have been colliding at a center of mass energy of \(\sqrt{s}=13.6\ TeV\), which is the highest energy ever reached in a collider and a little increase with respect to those of \(\sqrt{s}=13\ TeV\) of Run 2.  The increase in energy was accompanied by a rise in instantaneous luminosity, which has passed \(2\times10^{34}\ cm^{-2}s^{-1}\), bringing to almost 60 the number of proton-proton collisions per bunch crossing. This pile-up level led to a hard constraint on the triggers and detector sensitivity. Run 3 ended in 2026, with its last \(pp\) collisions in May and Pb-Pb collisions in June; then the LHC entered the Long Shutdown 3, in view of the High-Luminosity LHC. In 4 years of data-taking, the ATLAS experiment recorded \(332\ fb^{-1}\) \(pp\) collision data, double with respect to Run 2, and the biggest ever. This simulated data sample, together with the detector configuration described in Chapter 3, forms the basis of this work. The thesis is organised as follows. Chapter~2 outlines the Standard Model, with particular attention to the production and decay of the $Z$ boson into the muon final state on which the analysis is based. Chapter~3 introduces the ATLAS experiment, while Chapter~4 focuses on the Muon Spectrometer and the upgrades it has undergone. Chapter~5 describes the procedures adopted to reconstruct and identify the muons of interest, and Chapter~6 presents the analysis of the data and Monte Carlo samples. Chapter~7 summarises the results and draws the conclusions. "PER ADESSO"



\clearpage

% ==================================================================
%  CAPITOLO 2 - MODELLO STANDARD
% ==================================================================
\chapter{A Brief Overview of the Standard Model}


\lettrine[lines=3,loversize=0.1,lraise=0.05]{S}{} ince the 1930s, physicists have been trying to build a theory that connects Quantum Mechanics with Relativistic Mechanics. Numerous studies led to a brilliant insight, namely that the entire universe is made up of elementary entities, called particles, and that this 'Zoo' made up all the matter (and antimatter) we know. The most widely accepted theory with a wealth of experimental evidence is the Standard Model, formulated through Quantum Field Theory between the '60s and '70s. However, even though it provides an almost complete description of the entire cosmos, it is still incomplete. In fact, it doesn't include gravity, it doesn't explain dark matter, and it can't give us a measure for neutrino masses (and much more), and a whole thesis wouldn't be enough to describe everything. In this chapter, I'll go through the highlights of this theory, and the most well-known success that we probably all know, the evidence and discovery of the Higgs boson.

\section{Introduction}

The Universe as we know it is made up of two types of particles, distinguished by their spin value: fermions, with half-integer spin, which make up matter, and bosons, with integer spin, among which are the mediators of the fundamental interactions.
There are four known interactions: electromagnetic, strong, weak and gravitational, where the latter will not be discussed in this thesis. Fermions are divided into leptons and quarks, each with its own antiparticle. The leptons, reported in \cref{tab:sm_leptons}, are organized into three families, or
"generations", identical in quantum numbers and distinguished only by mass. The first generation constitutes all stable matter, while the other two decay towards
it. Charged leptons are subject to electromagnetic and weak interactions, 
while neutrinos interact only weakly. The electron is stable; the muon
it decays for weak interaction with a half-life of $2.2~\mathrm{\mu s}$, and the same goes for tau, but with a half-life of $2.9 \cdot 10^{-13}$~s. The quarks, shown in \cref{tab:quarks}, all have spin $1/2$ and fractional charge, $+2/3,|e|$ or $-1/3,|e|$. They show no internal structure, and current experimental limits place their size below $10^{-18}$~m. They exist in six flavors, distinguished by internal quantum numbers: up ($u$), down ($d$), charm ($c$), strange ($s$), bottom ($b$), and top ($t$). The $u$ and $d$ quarks are the lightest and have comparable masses, so much so that they are grouped into a
\begin{figure}[htpb]
    \centering
    \includegraphics[width=\textwidth]{TESI/particles.png}
    \caption{Schematic representation of the elementary particles of the Standard Model: quarks, leptons, gauge bosons and the Higgs boson. \cite{CERN-SM-chart}}
    \label{fig:particles}
\end{figure}
isospin doublets\footnote{Isospin was introduced by Heisenberg to describe the proton and neutron as two states of the same object with respect to the strong interaction; reflecting their similarity directly in the flavor symmetry between $u$ and $d$.}  $I = 1/2$, with third component $I_3 = +1/2$ for $u$ and $I_3 = -1/2$ for $d$. To the other flavors, a dedicated quantum number is associated: strangeness $S = -1$ for the quark $s$, charm $C = +1$ for the quark $c$,
beauty $B = -1$ for the quark $b$ e topness $T = +1$ for the quark $t$. Each quark has a corresponding antiquark with opposite electric charge and flavor numbers. Despite significant experimental efforts, free quarks have never been observed, as various studies suggest that they remain confined in bound states called hadrons. The intensity of each interaction is expressed by a coupling constant, which is a dimensionless parameter that measures its strength at a given vertex. However, even though they're called constants, they're not fixed at all, but change with the energy scale of the processes considered. In the case of
quantum chromodynamics (QCD), the constant $\alpha_s$ decreases as it increases
of the energy exchanged, i.e. as the distance between the quarks decreases. This means that at high energies $\alpha_s$ becomes small and quarks behave almost like free particles, a regime known as asymptotic freedom, which is what makes it possible to the perturbative calculation of high-energy processes at accelerators. At low energies, i.e. bigger distances, $\alpha_s$ grows until it is impossible to separate two quarks, which tells us that the energy needed to pull them apart produces new quark-antiquark pairs instead of releasing them, a regime known as confinement. Weak interactions violate quark flavor, and it’s precisely this property, described by the CKM matrix, that makes the CP violation discussed in Sec. 2.6 possible. Hadrons come in two main configurations: mesons, which are quark-antiquark bound states, and baryons, made up of three quarks (or three antiquarks, in the case of antibaryons); the proton, for example, is made of two up quarks and one down quark. In recent years, exotic states like tetraquarks and pentaquarks have also been observed, which don't change the picture described here. The reason why only certain combinations are possible is explained by QCD through the concept of color. In fact, every quark exists in three color states (red, green, and blue, with the corresponding anticolors for antiquarks) and observable hadrons are color singlets, meaning they are overall neutral. Color is a disctinct and independent property with respect to the flavor, and its the charge of the strong interactions. Now let's move on to bosons. These are the force mediators, meaning that particles interact by exchanging them with each other, and each type of interaction has its own characteristic mediator. The electromagnetic interaction is carried by the photon $\gamma$, which has no mass and no electric charge, giving it an infinite range. The strong interaction is mediated by eight gluons, which are also massless but carry color charge. This makes them self-interacting (unlike the gamma boson in the electromagnetic interaction), which sets QCD apart from QED and leads to the confinement of quarks. The weak interaction is instead mediated by bosons $W^\pm$ and $Z^0$, massive around $80$ GeV and $91$ GeV respectively. Their range is short precisely because of their mass, causing their interaction time to be much slower than that of strong or electromagnetic interactions. The $W^\pm$ bosons, carrying electric charge and are therefore the only mediators capable of change the flavor of fermions. Gravitational interaction is left out of this framework both because its hypothetical mediator, the graviton, has never been observed, and because its effects are so tiny that they can be ignored. Finally, to complete the particle content of the Standard Model there is the Higgs boson, the only scalar of the theory, with spin $0$ and mass of about $125$ GeV . Unlike gauge bosons, it does not mediate any interaction, but its role is to spontaneously break the electroweak symmetry, giving mass to the bosons $W^pm$ and $Z^0$ and, through Yukawa terms, to the fermions themselves. Predicted in the 1960s and observed at ATLAS and CMS in 2012, the mechanism that bears its name is discussed in detail in Sec. 2.7.
With these concepts just discussed, other questions that the SM can't explain are: why are they grouped into three families? And why in generations of mass?
\begin{table}[htpb]
    \centering
    \begin{tabular*}{\textwidth}{@{\extracolsep{\fill}}cccc@{}}
        \toprule
        \textbf{Generation} & \textbf{Doublet} & $\boldsymbol{Q}$ & \textbf{Mass} \\
        \midrule
        I &
        $\begin{pmatrix}\ \nu_e \\ e \end{pmatrix}_L$ &
        $\begin{matrix} 0 \\ -1 \end{matrix}$ &
        $\begin{matrix} < 0.8~\text{eV} \\ 0.511~\text{MeV} \end{matrix}$ \\
        \midrule
        II &
        $\begin{pmatrix}\ \nu_\mu \\ \mu \end{pmatrix}_L$ &
        $\begin{matrix} 0 \\ -1 \end{matrix}$ &
        $\begin{matrix} < 0.19~\text{MeV} \\ 105.7~\text{MeV} \end{matrix}$ \\
        \midrule
        III &
        $\begin{pmatrix}\ \nu_\tau \\ \tau \end{pmatrix}_L$ &
        $\begin{matrix} 0 \\ -1 \end{matrix}$ &
        $\begin{matrix} < 18.2~\text{MeV} \\ 1777~\text{MeV} \end{matrix}$ \\
        \bottomrule
    \end{tabular*}
    \caption{Lepton doublets of the Standard Model, with electric charge $Q$ in units of $e$.  Neutrino masses are upper limits.}
    \label{tab:sm_leptons}
\end{table}
\begin{table}[H]
    \centering
    \begin{tabular*}{\textwidth}{@{\extracolsep{\fill}}ccccc@{}}
        \toprule
        \textbf{Generation} & \textbf{Doublet} & $\boldsymbol{Q}$ & \textbf{Flavour} & \textbf{Mass} \\
        \midrule
        I &
        $\begin{pmatrix} u \\ d \end{pmatrix}_L$ &
        $\begin{matrix} +2/3 \\ -1/3 \end{matrix}$ &
        $\begin{matrix} I_3 = +1/2 \\ I_3 = -1/2 \end{matrix}$ &
        $\begin{matrix} 2.16~\text{MeV} \\ 4.67~\text{MeV} \end{matrix}$ \\
        \midrule
        II &
        $\begin{pmatrix} c \\ s \end{pmatrix}_L$ &
        $\begin{matrix} +2/3 \\ -1/3 \end{matrix}$ &
        $\begin{matrix} C = +1 \\ S = -1 \end{matrix}$ &
        $\begin{matrix} 1.27~\text{GeV} \\ 93.4~\text{MeV} \end{matrix}$ \\
        \midrule
        III &
        $\begin{pmatrix} t \\ b \end{pmatrix}_L$ &
        $\begin{matrix} +2/3 \\ -1/3 \end{matrix}$ &
        $\begin{matrix} T = +1 \\ B = -1 \end{matrix}$ &
        $\begin{matrix} 172.5~\text{GeV} \\ 4.18~\text{GeV} \end{matrix}$ \\
        \bottomrule
    \end{tabular*}
    \caption{Left-handed quark doublets of the Standard Model.In low-energy models of hadrons, quarks are considered dressed because of their interaction with gluons, which are the carriers of the strong force. These dressed quarks are called constituent quarks, and the mass in the table refers to the constituent quark.}
    \label{tab:quarks}
\end{table}




\section{Standard Model}

The Standard Model is a Gauge Theory \(SU(3)_C\ \otimes\ SU(2)_L\ \otimes\ U(1)_Y\) where \(SU(3)_C\) is the group of \(3\times3\) Special Unitary Matrices, associated with an internal colour quantum number. It represents Quantum Chromodynamics (QCD), the strong interaction, where both left-handed and right-handed quarks are subject to the same action of this group; that is, they are affected in a vector-like way, with only the colour component being rotated. The \(SU(2)_L\) is the group of \(2\times2\) Special Unitary Matrices, which acts only on left-handed particles and is therefore a chiral group. Unlike the group \(SU(3)_C\), it does not possess an internal quantum number; consequently, the associated Pauli matrices act by rotating the components of the family doublet. The \(U(1)_Y\), instead, is the Unitary Group of Hypercharge. The "1" indicates that the group has a single generator, corresponding to a phase. This is also a chiral group, since the left-handed hypercharge differs from the right-handed hypercharge. Associated with these groups is the Lagrangian of the Standard Model in the chiral gauge symmetry framework, expressed in the unbroken basis:
\begin{equation}
\mathcal{L}_{SM}\ \supset\ -\frac{1}{2}\mathrm{Tr}\left\{ G_{\mu\nu}(x)G^{\mu\nu}(x) \right\}-\frac{1}{2}\mathrm{Tr}\left\{ B_{\mu\nu}(x)B^{\mu\nu}(x) \right\}-\frac{1}{4}\ C_{\mu\nu}(x)C^{\mu\nu}(x)
\end{equation}
where the first term is a non-Abelian term associated with the group \(SU(3)_C\), describing eight generators:
\begin{equation}
G_\mu(x)=\sum_{A=1}^8 G_\mu^A \hat{T}^A
\end{equation}
where \(A=1,\dots,8\) labels the gluons, expressed through the Gell-Mann matrices. The second term is also a non-Abelian term, associated with the group \(SU(2)_L\) and describing three generators:
\begin{equation}
B_\mu(x)=\sum_{i=1}^3 B_\mu^i \frac{\hat{\sigma}^i}{2}
\end{equation}
where \(\hat{\sigma}^i\) are the Pauli matrices; the last term, instead, is an Abelian term associated with the group \(U(1)_Y\) and has just one generator. Therefore, in the unbroken basis of the SM, we have several vector bosons associated with gauge interactions: 8 vector bosons \(G_\mu\) associated with the strong interaction, and 4 vector bosons \(B_\mu^i\) and \(C_\mu\) associated with the electroweak theory. These gauge bosons are massless, and therefore possess only two degrees of freedom, corresponding to the two transverse polarisations with respect to the direction of motion. However, a model described entirely in terms of massless particles is in clear contradiction with experimental observations, as one can see by simply considering the \(W^{\mp}\) and \(Z^0\) bosons. The resolution of this problem had to await the 1960s, when Peter Higgs and Jeffrey Goldstone solved the problem of the spontaneous breaking of global and local gauge symmetries.

\section{Higgs Potential}
Using the Higgs doublet notation:
\begin{equation}
\phi = \begin{pmatrix} \phi^+ \\ \phi^0 \end{pmatrix}
\end{equation}
introduced to compensate for the quantum numbers of the \(SU(2)_L\) and \(U(1)_Y\) groups within the Yukawa terms, if we include the following terms in the SM Lagrangian:
\begin{equation}
\mathcal{L}_{SM} \supset -\phi^\dagger D_\mu D^\mu \phi +\lambda \phi^\dagger \phi(\phi\phi^\dagger-v^2)
\end{equation}
we can define the Higgs Potential:
\begin{equation}
V^H\left[ |\phi|^2\right] = \int d^4x\ \lambda|\phi|^2(|\phi|^2-v^2) \quad \text{with}\quad v^2=-\frac{\mu^2}{\lambda}>0
\end{equation}
where $\lambda$ is the Higgs self-coupling parameter and $v$ is the vacuum expectation value of the Higgs field. Moreover, $|\phi|^2$ is gauge invariant, and so $V^H$ is gauge invariant as well. Then, the potential follows:
\begin{equation}
V^H\sim|\phi|^4-|\phi|^2
\end{equation}
in full agreement with the renormalizability of the theory, that is, in the classical action there cannot exist operators with engineering dimension greater than the spacetime dimension.

\noindent The Higgs doublet has 4 degrees of freedom (i.e.\ 4 real numbers), that is the real part \(\mathrm{Re}(\phi^+,\phi^0)\) and the imaginary part \(\mathrm{Im}(\phi^+,\phi^0)\). However, it is convenient to switch to polar coordinates:
\begin{equation}
\phi = \begin{pmatrix} \phi^+ \\ \phi^0 \end{pmatrix} = e^{i\ \overrightarrow{\pi (x)}\ \frac{\overrightarrow{\sigma}}{2}} \begin{pmatrix} 0 \\ \frac{v+h(x)}{\sqrt{2}} \end{pmatrix}
\end{equation}
where \(h(x)\) is a fluctuation around the minimum and is a real number; indeed, for \(h(x)=0\), the minimum occurs at \(|\phi|=\pm v/\sqrt{2}\). The field \(\overrightarrow{\pi(x)}\) represents three symmetry rotations of the potential (three real numbers, therefore the number of degrees of freedom is conserved), such that \(V^H\) is invariant under these rotations. In particular, a valley of minima is formed, an infinite set of equivalent minima which, in quantum field theory, correspond to the minima of the Hamiltonian: that is, there are infinitely many vacuum states. These are connected by an \(SU(2)_L\) symmetry operation: the three rotations \(\overrightarrow{\pi(x)}\), which allow transitions from one vacuum state to another without any energy cost. The transition between different vacuum states gives rise to spontaneous symmetry breaking. The term ``breaking'' can however be misleading, because the full theory remains globally invariant under \(SU(2)_L\), and this phenomenon merely hides (rather than removes) the symmetry.

\section{Higgs Mechanism}
The Higgs Mechanism was introduced to explain how a local gauge symmetry can be broken. By fixing the unitary gauge:
\begin{equation}
e^{i\ \overrightarrow{\pi (x)}\ \frac{\overrightarrow{\sigma}}{2}} =1
\end{equation}
and due to locality, it is possible to remove the phase at every point in spacetime, i.e.\ to eliminate the dependence on \(\overrightarrow{\pi(x)}\). This would not have been possible in the case of a global symmetry, since there is no pointwise dependence on \(x\). In this way, the Higgs doublet reduces to (neglecting the \(h(x)\) term for simplicity):
\begin{equation}
\phi = \begin{pmatrix} 0 \\ v/\sqrt{2} \end{pmatrix}
\end{equation}
The kinetic term becomes:
\begin{equation}
-\phi^\dagger D_\mu D^\mu \phi \propto -\frac{1}{2} (0\ \ v) \left\{ \partial_\mu - ig_Y\frac{\hat{Y}}{2}C_\mu(x)-ig_wB^i_\mu(x)\frac{\hat{\sigma}^i}{2} \right\}^2 \begin{pmatrix} 0 \\ v \end{pmatrix}
\end{equation}
and we can neglect \(\partial_\mu\) because \(v=\mathrm{const}\). After a few simple steps, one finds that:
\begin{equation}
W^+=\frac{B^1_\mu-iB^2_\mu}{\sqrt{2}}=\left( W^-_\mu \right)^\dagger \quad, \ Z^0_\mu=- \frac{g_YC_\mu-g_wB^2_\mu}{g} \quad, \ g=\sqrt{g_Y^2+g_w^2}
\end{equation}
where $g_{Y}$ and $g_w$ the Hyperchaged and Weak coupling constant respectivetly. Now we can write Eq. 2.11 as:
\begin{equation}
-\phi^\dagger D_\mu D^\mu \phi \propto\frac{v^2}{2} \left\{ g_w^2 \frac{W_\mu^+W_\mu^-}{2} + g^2\frac{Z^0_\mu Z^0_\mu}{4} \right\}
\end{equation}
where \(\frac{W_\mu^+W_\mu^-}{2}\) is a complex field of the type \(\phi^\dagger \phi\) and \(\frac{Z^0_\mu Z^0_\mu}{4}\) is a real field of the type \(|\phi|^2\). This yields the masses of the weak bosons:
\begin{equation}
m_W^2=v^2\frac{g_w^2}{4} \quad, \ m_Z^2=v^2\frac{g^2}{4}
\end{equation} 
So the Higgs mechanism gave mass to 3 out of 4 bosons of the weak interactions, with mass proportional to the coupling constants and to the vacuum expectation value of the Higgs field $v$. These masses are also related through the Weinberg angle:
\begin{equation}
\frac{m_W^2}{m_Z^2}=\cos^2\theta_w\ \quad\ , \ \frac{g_Y^2}{g^2}=\sin^2\theta_w
\end{equation}
from which one experimentally obtains that:
\begin{equation}
m_W\sim 80\ \text{GeV}\ \quad\ , \ m_Z\sim 91\ \text{GeV}
\end{equation}
With this Weinberg angle, it is also possible to relate:
\begin{equation}
\begin{pmatrix} A_\mu \\ Z^0_\mu \end{pmatrix} = \begin{pmatrix} \sin\theta_w & \cos\theta_w \\ \cos\theta_w & -\sin\theta_w \end{pmatrix} \begin{pmatrix} B_\mu^2 \\ C_\mu \end{pmatrix}
\end{equation}
where \(A_\mu\) is the photon, expressed as the combination orthogonal to \(Z^0_\mu\); it remains massless. From the covariant derivative of the quarks, considering only \(B^2_\mu \frac{\hat{\sigma}^3}{2}\), and isolating \(A_\mu\), one finds that the electric charge is:
\begin{equation}
D_\mu q\propto-ig\ \sin\theta_w\ \cos\theta_w\ A_\mu\ \left( \frac{\hat{Y}}{2} +\frac{\hat{\sigma}^3}{2}\frac{1-\gamma^5}{2} \right) = -i\ e\ A_\mu\ \hat{Q}_e
\end{equation}
\footnote{$\mathcal{P}_L=\frac{1-\gamma^5}{2}$  is the Left-chirality projection operator, while $\mathcal{P}_R=\frac{1+\gamma^5}{2}$ is the Right-chirality projection operator}In the Higgs Mechanism, when a vector boson acquires mass, it gains an additional degree of freedom: the longitudinal polarisation. In our case, three such degrees are required, one for each boson that has acquired mass, and it is natural to ask where these extra degrees of freedom come from. They are the three Goldstone bosons \(\overrightarrow{\pi(x)}\), which, upon reorganisation of the basis, reappear as longitudinal polarisations:
\[\ll\ "W^\mp_\mu\ \text{and}\ Z^0_\mu\ \text{have eaten the Goldstone bosons and got mass.}"\ \gg \]
This is the so-called ``broken basis'' of the Standard Model.

\section{Yukawa Term}
Using the expression of the Higgs doublet in the unitary gauge, let us see what happens if we diagonalise the Yukawa terms:
\begin{equation}
\mathcal{L}_{SM} \supset -\overline{q}_f \left\{ U_{fg}\phi^c\frac{1+\sigma^3}{2}+D_{fg}\phi\frac{1-\sigma^3}{2} \right\}\frac{1+\gamma^5}{2}q_g+h.c.
\end{equation}
where \(\phi^c=i\sigma^2\phi^*\) is the charge conjugate and $U_{fg}$, $D_{fg}$ are the Yukawa couplings and they are different between the different families . Collecting all the families into a single large vector \(\overline{q}=(\overline{u},\overline{c},\overline{t},\overline{d},\overline{s},\overline{b})\), explicitly writing the Hermitian conjugate and performing a chiral rotation, we obtain a diagonalised version:
\begin{equation}
\overline{q}=\overline{X}\left\{ \begin{pmatrix} A_U^\dagger & 0 \\ 0 & A_D^\dagger \end{pmatrix} \frac{1+\gamma^5}{2} + \begin{pmatrix} B_U^\dagger & 0 \\ 0 & B_D^\dagger \end{pmatrix} \frac{1-\gamma^5}{2}  \right\}
\end{equation}
\begin{equation}
q=\left\{ \begin{pmatrix} A_U & 0 \\ 0 & A_D \end{pmatrix} \frac{1-\gamma^5}{2} + \begin{pmatrix} B_U & 0 \\ 0 & B_D \end{pmatrix} \frac{1+\gamma^5}{2}  \right\}X
\end{equation}
\begin{equation}
\mathcal{L}_{SM} \supset\frac{v}{\sqrt{2}}\overline{X} \left\{ \begin{pmatrix} A_U^\dagger UB_U & 0 \\ 0 & A_D^\dagger DB_D \end{pmatrix} \frac{1+\gamma^5}{2} + \begin{pmatrix} B_U^\dagger U^\dagger A_U & 0 \\ 0 & B_D^\dagger D^\dagger A_D \end{pmatrix} \frac{1-\gamma^5}{2}  \right\}X
\end{equation}
where the matrices \(A\) and \(B\) are unitary. By the fundamental theorem of algebra, any matrix can always be diagonalised via a biunitary transformation, i.e.\ with real eigenvalues. Therefore, with a suitable choice of \(A\) and \(B\), one can write:
\begin{equation}
\begin{pmatrix} m_u & 0 & 0 \\ 0 & m_c & 0 \\ 0 & 0 & m_t \end{pmatrix} = \frac{v}{\sqrt{2}}\ A^\dagger_U UB_U = \frac{v}{\sqrt{2}} B^\dagger_U U^\dagger A_U
\end{equation}
\begin{equation}
\begin{pmatrix} m_d & 0 & 0 \\ 0 & m_s & 0 \\ 0 & 0 & m_b \end{pmatrix} = \frac{v}{\sqrt{2}}\ A^\dagger_D DB_D = \frac{v}{\sqrt{2}} B^\dagger_D D^\dagger A_D
\end{equation}
So, all the masses in the Standard Model are proportional to \(v\), and the Yukawa terms reduce to a simple sum over the quark masses:
\begin{equation}
\mathcal{L}_{SM} \supset -\overline{q}_f \left\{ U_{fg}\phi^c\frac{1+\sigma^3}{2}+D_{fg}\phi\frac{1-\sigma^3}{2} \right\}\frac{1+\gamma^5}{2}q_g\ + h.c. =\sum_f m_f\overline{X}_fX_f
\end{equation}
By also including \(h(x)\), we can obtain the Higgs-Quark interaction vertex, which are proportional to the quark masses, as shown in \cref{fig:hver}:
\begin{figure}[H]
    \centering
    \adjustbox{valign=c}{%
    \(\overline{X}_f\, h(x)\, X_f\)
    }
    \hspace{0.5cm}
    \adjustbox{valign=c}{%
    \includegraphics[width=0.25\linewidth]{TESI/hver.png}
    }
    \caption{Higgs-Quark interaction vertex, shown as a Feynman diagram.}
    \label{fig:hver}
\end{figure}
\noindent The change of basis induced by the chiral rotation also affects the kinetic term of the quarks, in particular one term of the weak interactions, giving rise to charged and neutral current interactions:
\begin{equation}
\begin{aligned}
\overline{q}\gamma^\mu D_\mu q ={}& \overline{X}\gamma^\mu\Bigg\{\partial_\mu
  - i(G_\mu,A_\mu,Z^0_\mu) \\[0.4cm]
&\quad - ig_w\frac{\hat{V}_{CKM}W^+_\mu\sigma^+ + (\hat{V}_{CKM})^\dagger W^-_\mu\sigma^-}{\sqrt{2}}\,
  \frac{1-\gamma^5}{2}\Bigg\}X
\end{aligned}
\end{equation}
where \(- i(G_\mu,A_\mu,Z^0_\mu)\) encloses all the neutral current terms, which do not mix the flavour, whereas the second term represents the charged current, which can mix the flavour. The presence of \(\hat{V}_{CKM}\) implies that in weak interactions the charged current interactions mix flavour, since the \(\hat{V}_{CKM}\) matrix gives the probability of mixing between the different quark families (the mixing is larger on the diagonal, especially between the first and second family). The unitarity of \(\hat{V}_{CKM}\), i.e.\ \(\hat{V}_{CKM}(\hat{V}_{CKM})^\dagger=A^\dagger_U A_D A^\dagger_D A_U=\mathbb{I}\), together with its complex phase, gives rise to CP symmetry violation.

\section{CP Violation}
\label{cp_violation}
Weak interactions violate parity (P), charge conjugation (C) and CP symmetry. The latter is the product of the first two and plays a fundamental role in explaining the asymmetry between matter and antimatter in the Universe, although not with sufficient magnitude. Performing a CP transformation on the charged current weak interaction term, we get:
\begin{equation}
\begin{aligned}
\mathcal{L}_{SM} &\supset
  \ccterm{\hat{V}_{CKM}W^+_\mu\sigma^+ + (\hat{V}_{CKM})^\dagger W^-_\mu\sigma^-} \\[0.5cm]
&\overset{CP}{\longrightarrow}\
  \ccterm{(\hat{V}_{CKM})^*W^+_\mu\sigma^+ + (\hat{V}_{CKM})^\dagger W^-_\mu\sigma^-}
\end{aligned}
\end{equation}
The only thing that prevents this term from being invariant is whether:
\begin{equation}
(\hat{V}_{CKM})^*\overset{?}{=}\hat{V}_{CKM}
\end{equation}
If the matrix is real, the term is invariant under CP; otherwise, it is not.
We know that \(\hat{V}_{CKM}\) couples up-type quarks to down-type quarks in charged current weak interactions:
\begin{equation}
\hat{V}_{CKM}=\begin{pmatrix} \hat{V}_{ud} & \hat{V}_{us} & \hat{V}_{ub} \\ \hat{V}_{cd} & \hat{V}_{cs} & \hat{V}_{cb} \\ \hat{V}_{td} & \hat{V}_{ts} & \hat{V}_{tb} \end{pmatrix}
\end{equation}
In the most general case, it contains \(n^2\) complex entries, where \(n\) is the number of families. We may then ask: if it were real, how many angles would it contain? A real \(n\times n\) matrix has \(\#\text{angles}=\frac{n(n-1)}{2}\) and \(\#\text{phases}=n^2-\#\text{angles}=\frac{n(n+1)}{2}\). For \(n=3\), there are three rotation angles and six phases. However, not all of these six phases are physical. By performing a phase transformation of the form:
\begin{equation}
q_u=e^{i\hat{\phi}_u}q_u \quad q_d=e^{i\hat{\phi}_d}q_d
\end{equation}
this affects all the terms in the Lagrangian; however, they are diagonal and therefore remain invariant. The only term that is not invariant is:
\begin{equation}
\hat{V}_{CKM}\longrightarrow e^{-i\hat{\phi}_u}\hat{V}_{CKM}e^{i\hat{\phi}_d}
\end{equation}
This allows us to absorb \(2n-1\) unphysical phases, and therefore the effective number of physical phases is given by:
\begin{equation}
\#\text{Physical phases}= \#\text{Phases} - \#\text{Unphysical phases}= \frac{(n-1)(n-2)}{2}
\end{equation}
Then, for \(n=3\), \(\hat{V}_{CKM}\) is parametrised by 3 rotation angles and 1 complex phase. 
This shows that a physical phase can exist and can't be removed by redefining the fields, but at the same time it doesn't prove that this phase has to be different from zero. Its non-zero value is an experimental fact.
Therefore the Cabibbo--Kobayashi--Maskawa matrix cannot be real, and the presence of this single phase leads to the breaking of CP symmetry in the Standard Model.

\section{The Higgs Boson}

Restoring the fluctuation $h(x)$ in the doublet of Eq. 2.10,
$\phi = \left(0,\ (v+h)/\sqrt{2}\right)^T$, the potential becomes
$V(h) = \lambda v^2 h^2 + \lambda v\, h^3 + \frac{\lambda}{4} h^4$. Comparing the quadratic term with the mass term of a real scalar field, $\frac{1}{2} m_H^2 h^2$, one obtains
\begin{equation}
    m_H = \sqrt{2\lambda}\, v ,
    \label{eq:higgs_mass}
\end{equation}
with $v = (\sqrt{2}\,G_F)^{-1/2} \simeq 246$~GeV and $G_F$ the Fermi constant. Since the self-coupling $\lambda$ is presently unknown, the value of the SM Higgs boson mass $m_H$ cannot be predicted. However, theoretical  arguments place approximate upper and lower bounds to $m_H$. At hadron colliders the Higgs boson is produced mainly through gluon fusion ($gg \to H$, via a top-quark loop), followed by vector boson fusion
($qq \to qqH$) and associated production with a $W$, a $Z$ or a $t\bar{t}$ pair. For a light Higgs boson the dominant decay is $H \to b\bar{b}$, but the
discovery channels are $H \to \gamma\gamma$ and $H \to ZZ^* \to 4\ell$ \footnote{$Z*$ it's not the real Z boson, but virtual, meaning it exists only within Heisenberg's uncertainty principle}, also known as "the golden channel", which have a small branching ratio but an excellent mass resolution representing the most clean decay channel for the Higgs boson. Before the LHC, LEP set the lower bound $m_H > 114.4$~GeV at 95\% confidence level, while electroweak precision fits indicated $m_H < 152$~GeV. In 2011 the Tevatron and the LHC excluded most of the remaining range, leaving a window around $125$~GeV. On 4 July 2012 ATLAS and CMS announced its observation, with significances of $5\sigma$ for a mass of $m_H = 125.3 \pm 0.6$~GeV (CMS) \cite{PDG2012}.

\begin{figure}[htpb]
    \centering
    \includegraphics[width=\textwidth]{TESI/golden.png}
    \caption{Expected and observed four-lepton invariant mass distribution $m_{4l}$ for all Higgs boson candidates.The uncertainty in the predictions is shown by the hatched band, which includes the theoretical uncertainties of the SM cross-section for the signal and the background processes. The observed number of events and the expected signal and background yields in the mass range $115\ \mathrm{GeV}<m_{4l}<130\ \mathrm{GeV}$}.
    \label{fig:placeholder}
\end{figure}



\chapter{Charged particles in the Standard Model}

In the previous chapter, we saw how the Standard Model describes elementary particles and their interactions, and how the Higgs boson has been observed in final states like the so-called golden channel, $H \to ZZ^* \to 4\ell$ (with $\ell = e, \mu$). The nickname "Golden", as already noted in the previous section, comes from the fact that among all its possible final states, although it is not the most common, it is the cleanest, especially if the four leptons are muons. This is because, more generally, final states with only muons allow the signal to be isolated more easily, unlike those that contain gluon jets or many hadrons which, decaying in turn, produce a much higher background. This is due to the properties of the muon, which being a lepton about 207 times heavier than the electron, is not affected by the strong interaction and, with an average lifetime of about $2.2~\mu\text{s}$, lives long enough to pass through the entire detector before decaying, managing to be identified and measured precisely. In this chapter we will describe the passage of muons through matter, that is, the energy loss by ionization given by the Bethe-Bloch formula and multiple Coulomb scattering, along with the measurement of momentum through the track's sagitta in a magnetic field. These elements will lead to the general parametrization of the momentum resolution, whose terms will be measured in Chapter 6, and will form the basis for the description of the ATLAS detector and its muon spectrometer in Chapters 3 and 4.

\section{Bethe-Block formula}

A charged particle passing through a material loses energy mainly in collisions with atomic electrons, which get excited or ionized. For charged particles heavier than the electron, the average energy loss per unit length of the material is described by the Bethe-Bloch (BB) formula \cite{PDG2026}. For electrons, the treatment is different, both because of the kinematics and because for them radiative losses already dominate above just a few tens of MeV. 
\begin{equation}
  -\left\langle \frac{\mathrm{d}E}{\mathrm{d}x} \right\rangle
  = K z^2 \, \frac{Z}{A} \, \frac{1}{\beta^2}
  \left[ \frac{1}{2} \ln \frac{2 m_e c^2 \beta^2 \gamma^2 W_{\max}}{I^2}
  - \beta^2 - \frac{\delta(\beta\gamma)}{2} - \frac{C}{Z}
  + \frac{1}{8} \left( \frac{W_{\max}}{\gamma M c^2} \right)^{2} \right]
  \label{eq:bethe_muon}
\end{equation}

\begin{equation}
  W_{\max} = \frac{2 m_e c^2 \beta^2 \gamma^2}
  {1 + 2\gamma \, m_e/M + \left( m_e/M \right)^2}
  \label{eq:wmax}
\end{equation}

\noindent The equation \cref{eq:bethe_muon} gives the average value\footnote{Energy loss is the result of a finite number of collisions, and the actual loss fluctuates from event to event around this value.} of the energy loss due to ionization and excitation of a particle per unit length of material crossed, the trend of which is shown in \cref{fig:stopping} for a $\mu^+$ in Copper. In these equations $K = 4\pi N_A r_e^2 m_e c^2 \approx 0.307\ MeV mol^{-1} cm^2$, while $z=\pm 1$ it's the charge of the incoming particle in units of e. Z and A are the atomic number and atomic mass (in g/mol) of the material, $m_e c^2$ it's the electron's rest energy, M is the mass of the incoming particle, $\beta=v/c$ is the relativistic speed of the incoming particle, and $\gamma$ is the Lorentz factor. I is the average excitation energy of the material (332 eV for Copper) while $\delta(\beta\gamma)$ is the density correction and $C/Z$ the shell correction, discussed later. The last term in the parentheses is the correction for spin-1/2 particles, which only becomes relevant at high energies. Finally, $W_{max}$ is the maximum energy that the particle can transfer to an electron in a single collision. 
\begin{figure}[htpb]
    \centering
    \includegraphics[width=\textwidth]{stopping_power_mu_Cu.pdf}
    \caption{Trend of the stopping power $\langle -\mathrm{d}E/\mathrm{d}x\rangle$ for $\mu^+$ in Copper as function of $\beta\gamma = p/Mc$ (upper axis) and the muon momentum (lower axis). The vertical bands separate three different regimes that require different adjustments, like density correction or shell correction. In the interval $0.001\leq\beta\gamma\leq0.1$ we have the maximum stopping power where the incoming particle is complitely absorbed. As $\beta\gamma$ increases, the minimum of ionization of the incident particle characterized by a red dott-line is reached. The green dashed line shows the BB formula without density correction. The orange dotted line shows the radiative losses. $E_{\mu c}$ is the critical energy of the medium where the energy loss due to ionization equals that due to radiation, and it depends on the type of material.}
    \label{fig:stopping}
\end{figure}
\noindent On the y-axis, the braking power is shown, that is, the energy loss per unit length divided by the material's density. The BB formula depends on the material mainly through $Z/A$ and, in a logarithmic way, through $I$: expressed in this way, the magnitude varies little from one material to another. To get the loss per unit length, we just need to multiply by the density of the material (for Copper $\rho = 8.96~\mathrm{g/cm^3}$). On the x-axis are shown $\beta\gamma = p/Mc$ (upper axis) and the muon momentum (lower axis). The vertical bars separate the regimes in which different approximations apply. For $\beta\gamma \lesssim 0.005$ the Lindhard-Scharff regime applies, where the incoming particle is slower than the electrons in the outermost shell of the atoms in the material, and the stopping power is proportional to $\beta$. For $0.005 \lesssim \beta\gamma \lesssim 0.05$ instead, the Anderson and Ziegler parameterizations are used, because the particle's speed is comparable to that of the electrons in the outermost shell of the atoms in the medium. In both regions, the stopping power is very high and the muon comes to a stop in a very small thickness of material. The BB formula couldn't be applied to these two regimes, so the shell correction $-C/Z$ was introduced, which allows the formula to be used up to $\beta\gamma \approx 0.05$. In the interval $0.1 \lesssim \beta\gamma \lesssim 1$ factor:
\begin{equation}
  -\frac{\mathrm{d}E}{\mathrm{d}x} \propto \frac{1}{\beta^2}
\end{equation}
dominates and the energy loss decreases quickly as the momentum increases, until it reaches a minimum for $\beta\gamma \approx 3$--$4$. Particles with this energy are called minimum ionizing particles (MIP), where the particle keeps losing energy to the medium, but at the lowest possible rate (as shown in \cref{fig:stopping}, it goes around $1\text{-}2\ \mathrm{MeV\ g/cm^2}$). Beyond the minimum, the logarithmic term slowly increases the stopping power (relativistic rise):
\begin{equation}
  -\frac{\mathrm{d}E}{\mathrm{d}x} \propto \ln(\beta\gamma)
\end{equation}
both because of the increase in $W_{max}$ and because, as $\gamma$ increases, the electric field of the particle begins to polarize the electric field of the atoms in the medium. This way, a density correction term is introduced into the formula $\delta(\beta\gamma)$. At even higher energies, radiative processes come into play: bremsstrahlung, $e^+e^-$ pair production, and photonuclear interactions. Radiative loss increases approximately linearly with energy and, beyond the critical energy $E_{\mu c}$, i.e. the energy at which radiative losses and losses by ionization equal, it becomes dominant. Ionization losses don't stop, but they become an ever smaller fraction of the total. These features explain why the muon is so penetrating. The bremsstrahlung cross-section scales like $1/M^2$, so for a muon it is suppressed by a factor $(m_e/m_\mu)^2 \approx 2\times10^{-5}$ compared to an electron: in copper the electron's critical energy is about $20~\mathrm{MeV}$, that of the muon of a few hundred GeV. In the energy range of interest, the muon therefore doesn’t produce electromagnetic showers and, since it doesn’t interact strongly, neither hadronic showers: it loses energy almost only through ionization, at a rate close to the minimum, and behaves like a MIP. It can thus pass through large thicknesses of material and be identified and measured in detectors placed after the absorbing material (calorimeters). In the next Sections, will describes how its momentum is obtained from the curvature of the track under a magnetic field, whose integral along the trajectory,  $\int B\,\mathrm{d}l$, is called bending power: basically, we will measure the sagitta. The precision of this measurement is limited by passing through the material in two ways. The first is energy loss fluctuations, which actually result from a finite number of collisions, and sometimes a single collision gives a lot more energy to one electron. The second is multiple Coulomb scattering off the nuclei of the atoms in the medium, which deflects its trajectory. At high momentum, on the other hand, the sagitta becomes small and what really matters is how accurately we can measure the positions.

\section{Multiple Coulomb scattering}
A charged particle passing through a material is deflected by a large number of small-angle collisions, mainly due to Coulomb scattering on the nuclei. By the central limit theorem, the sum of these random deflections produces an almost Gaussian angular distribution\cite{PDG2026}. Projecting the deflection onto a plane that contains the initial direction, $98\%$ of the central part of the distribution is well approximated by a Gaussian of width
\begin{equation}
  \theta_0 = \frac{13.6~\mathrm{MeV}}{\beta c p}\, z
  \sqrt{\frac{x}{X_0}}
  \left[ 1 + 0.038 \ln\frac{x z^2}{X_0 \beta^2} \right],
  \label{eq:theta0}
\end{equation}
where $x$ is the thickness crossed, $z$ is the charge of the particle ($z=1$ for the muon), and $X_0$ is the radiation length of the material, which describe the length of a material type needed to reduce a particle's initial energy by 1/e. If the material has multiple layers, just use Eq. 3.5 once for the total thickness in radiation lengths, rather than adding up the contributions of each layer in quadrature. Moreover, $X_0$ decreases as $Z$ (the atomic number) increases: for the same thickness, lighter materials bend particles less.   

\section{Momentum measurement: the sagitta}
In a uniform magnetic field $B$, a particle with charge $ze$ moves along a helix. By equating the Lorentz force to the centripetal force, the component of momentum perpendicular to the field is related to the curvature radius $R$ of the trajectory by\cite[Sez.~35.12]{PDG2026}
\begin{equation}
  p\,[\mathrm{GeV}/c] = 0.3\, z\, B\,[\mathrm{T}]\; R\,[\mathrm{m}] .
  \label{eq:pBR}
\end{equation}
To determine $R$, and then $p$, we measure the position of the track
at several points. In the simplest case of three points, at the beginning, in the middle and at the end of a track segment of length $L$, the sagitta $s$ is the distance
of the central point from the line that connects the ends. For a slightly curved track ($L \ll R$) we have
\begin{equation}
  s = R\left(1 - \cos\frac{\alpha}{2}\right) \simeq \frac{L^2}{8R}
  = \frac{0.3\, z B L^2}{8\, p},
  \label{eq:sagitta}
\end{equation}
where $\alpha \simeq L/R$ is the deflection angle. The sagitta is inversely proportional to the momentum: a high-momentum track is almost straight. Eq. 3.7 it can be written as $p = k/s$, with
$k = 0.3\,zBL^2/8$ known. By propagating the error, we get
$\sigma_p = (k/s^2)\,\sigma_s = (p/s)\,\sigma_s$, that is
\begin{equation}
  \frac{\sigma_p}{p} = \frac{\sigma_s}{s} .
  \label{eq:sigmap_sigmas}
\end{equation}
The relative error on the momentum is the same as that on the sagitta: to estimate the resolution, it's enough to understand what makes the measurement of $s$ uncertain. There are three main causes, which we will look at one by one in the next section.

\section{Spatial, Multiple Scattering and Energy fluctuation resolution}
If each point is measured with resolution $\sigma_x$, the sagitta
$s = x_2 - (x_1 + x_3)/2$ has an error
$\sigma_s = \sqrt{\sigma_x^2 + \sigma_x^2/4 + \sigma_x^2/4}
= \sqrt{3/2}\,\sigma_x$. Dividing by Eq. 3.7 we get
\begin{equation}
  \left.\frac{\sigma_p}{p}\right|_{\mathrm{pos}}
  = \sqrt{\frac{3}{2}}\,\frac{8\,\sigma_x\, p}{0.3\, z B L^2} .
  \label{eq:res_pos3}
\end{equation}
With $N \gtrsim 10$ evenly spaced points the coefficient changes, but the dependence stays the same~\cite[Sez.~35.12]{PDG2026}:
\begin{equation}
  \left.\frac{\sigma_p}{p}\right|_{\mathrm{pos}}
  = \frac{\sigma_x\, p}{0.3\, z B L^2}\sqrt{\frac{720}{N+4}} .
  \label{eq:res_pos}
\end{equation}
Since sagitta decreases as $1/p$ while $\sigma_s$ remains
constant, this contribution grows linearly with the momentum. It reduces as $1/(BL^2)$: increasing the lever arm is way more effective than increasing the magnetic field. Even with perfectly measured points, multiple scattering randomly deflects the track and changes its sagitta. After a path length $L$ of material, the central point's deviation from the line joining the ends has a root mean square (RMS) value~\cite[Sez.~34.3]{PDG2026}
\begin{equation}
  s_{\mathrm{MS}} = \frac{L\,\theta_0}{4\sqrt{3}} ,
  \label{eq:s_ms}
\end{equation}
with $\theta_0$ given by Eq. 3.5. This fake sagitta acts as an error on the measurement $\sigma_s = s_{\mathrm{MS}}$. For a muon
($z = 1$), using Eq. 3.8 and neglecting the logarithmic term, i.e. with
$\theta_0 \simeq \frac{13.6~\mathrm{MeV}}{\beta c p}\sqrt{L/X_0}$,
we get:
\begin{equation}
  \left.\frac{\sigma_p}{p}\right|_{\mathrm{MS}}
  = \frac{s_{\mathrm{MS}}}{s}
  = \frac{2\,p\,\theta_0}{\sqrt{3}\cdot 0.3\, B L}
  \simeq \frac{0.05}{\beta\, B \sqrt{L X_0}} ,
  \label{eq:res_ms}
\end{equation}
with $p$ in GeV/$c$, $B$ in tesla and $L$, $X_0$ in meters. The deviation due to scattering increases as the momentum decreases, but in the same way, the true sagitta also increases: their ratio doesn't depend on $p$, so this contribution is constant. It is reduced by increasing the field and decreasing the material along the track. Before being measured, the particle passes through material in which
loses energy (Sec. 3.1). The average loss is
corrected for, but its fluctuations $\sigma_{\Delta E}$ from event to event remain as uncertainty on the momentum. For a relativistic muon
$E \simeq pc$, so $\sigma_p \simeq \sigma_{\Delta E}/c$ and
\begin{equation}
  \left.\frac{\sigma_p}{p}\right|_{\Delta E}
  \simeq \frac{\sigma_{\Delta E}}{p\,c} .
  \label{eq:res_eloss}
\end{equation}
Since beyond the minimum of the BB formula $\sigma_{\Delta E}$ it depends less by the momentum, this term decreases as
$1/p$~\cite{ATLAS-MCP-2023}. The three contributions are independent and they add up in quadrature:
\begin{equation}
  \frac{\sigma_p}{p} = \frac{a}{p} \oplus b \oplus c\,p ,
  \label{eq:res_tot}
\end{equation}
where $a \simeq \sigma_{\Delta E}/c$ comes from fluctuations in the
energy loss~\eqref{eq:res_eloss}, $b$ from the multiple
scattering~\eqref{eq:res_ms} and $c$ from the spatial resolution of the
points~\eqref{eq:res_pos}. At low momentum the first two terms dominate, at high momentum the third one does. So, to measure high-momentum muons well, we need a strong field, a long lever arm, little material along the track, and detectors with excellent abd a good 

% ==================================================================
%  CAPITOLO 3 - IL RIVELATORE ATLAS
% ==================================================================
\chapter{The ATLAS Detector}

\lettrine[lines=3,loversize=0.1,lraise=0.05]{T}{} he Large Hadron Collider (LHC) is the biggest and most powerful particle accelerator ever built. It’s at CERN, near Geneva, inside a circular tunnel about 27 km around. Almost 100 meters underground, it collides together proton beams (and sometimes heavy ions) at super high energies, letting scientists check out new elementary particles and the fundamental forces that show up from the collisions. There are four main experiments around the ring: ATLAS and CMS (general-purpose detectors), LHCb (for beauty quark physics), and ALICE (looking at heavy ions and quark-gluon plasma). This chapter explains how the various detectors that make up the ATLAS experiment are arranged, going through the reference system adopted by the ATLAS collaboration and focusing the discussion on the muon spectrometer (see Chapter 4), an essential element for analyzing its performance and muon resolution during Run 3. Also, a brief about the ATLAS trigger system will be included in this chapter. 

\section{ATLAS coordinate system}

The reference system used by the ATLAS collaboration places the origin of the axes at the geometric center of the ATLAS detector, that is, at the nominal interaction point (IP). A schematic representation of the axes is shown in \cref{fig:CS}, where the z-axis is oriented along the beam pipe, the y-axis points upwards, and the x-axis runs from the IP to the center of the LHC ring. In the (x,y) plane transverse to the beam, cylindrical coordinates (r, \(\phi\)) are used, where \(\phi\subset [0,2\pi]\) is the azimuthal angle between the x-axis and the particle's transverse momentum. The polar angle \(\theta \subset [0,\pi]\), on the other hand, is the angle between the z-axis and the particle's momentum. Therefore, \(\phi\) helps to understand the particle's position around the z-axis, while theta tells how far it is from it. Although rotations in \(\theta\) around the z-axis are invariant, since ATLAS has cylindrical symmetry, a different coordinate called pseudorapidity \(\eta = -\ ln \left[ tg\left( \theta/2 \right) \right] \) is used; see \cref{tab:eta-theta}. The physical meaning stays the same, but the difference is that for particles coming from the same collision, their cones have different \(\Delta\theta\) but the same \(\Delta\eta\). With this, we can also check out the angular distance \(\Delta R = \sqrt{\Delta\eta^2 + \Delta\phi^2}\).

\begin{figure}[H]
    \centering
    \includegraphics[width=\linewidth]{TESI/CS.jpeg}
    \caption{Schematic representations of ATLAS's coordinate system}
    \label{fig:CS}
\end{figure}
\noindent
\begin{table}[H]
  \centering
  \begin{tabular}{lccccc}
    \toprule
    \(\theta\) & \(90^\circ\) & \(40^\circ\) & \(15^\circ\) & \(5.7^\circ\) & \(0^\circ\)\\
               & \footnotesize(y-axis) & & & & \footnotesize(z-axis) \\
    \midrule
    \(\eta\)   & \(0\) & \(1\) & \(2\) & \(3\) & \(\infty\) \\
    \bottomrule
  \end{tabular}
    \caption{Schematic conversion of the angle \(\theta\) into pseudorapidity \(\eta\)}
    \label{tab:eta-theta}
\end{table}

\section{ATLAS layout during Run 3}

The detector is 25 m high and 44 m long; its overall weight is approximately 7000 t. Based on the coordinate system we just introduced, let's try to understand how the various ATLAS detectors are distributed. A quick overview of the Run 3 configuration is shown in \Cref{fig:ATLASexp}. Moving outwards from the beam line, within the 2 T xial magnetic field of the solenoid, we have: the Inner Detector system (ID) that allow us to reconstruct the tracks of the charged particles within \(|\eta|<2.5\).this is composed by the High-Granularity silicon pixel detector that covers the vertex region providing four measurements per track, where generally the first hit start in the Insertable B-Layer (IBL) \cite{ATLAScollab2, ATLASibl}, outside them we have the Semiconductor Tracker (SCT), a silicon microstrip tracker that provides eight measurements per track. Moreover, these silicon detectors are integrated with Transition Radiation Trackers (TRTs) that provide information on the electronic component within \(|\eta|=2.0\). Outside the solenoid, the calorimeter system gives full solid-angle coverage, going up to \(|\eta|<4.9\). The electromagnetic calorimeters sit in the \(|\eta|<3.2\) region, especially in the barrel and endcap, and they’re high-granularity, made of lead or Liquid Argon (LAr). As for the hadronic calorimeters, three are in the \(|\eta|<1.7\) area (barrel), made of alternating steel and scintillator layers, and two are in the \(1.5<|\eta|<3.2\) region (endcap), made of alternating copper and LAr layers. To cover the rest of the solid angle, the forward direction has smaller calorimeters in the same setup. The Muon Spectrometer (MS) comprises separate trigger and high-precision tracking chambers in order to measure the bending of muon tracks in the field produced by the superconducting air-core toroidal magnets, whose field integral lies between 2.0 Tm and 6.0 Tm over most of the detector. Then, we have three stations of precision chambers, where the innermost one of the end-caps, within \(|\eta|>1.3\), has been replaced by the New Small Wheel (NSW) detectors for Run 3. The other two stations, instead, each consist of layers of Monitored Drift Tubes (MDTs), which allow us to reconstruct tracks in the barrel. Installed during the Long Shutdown~2 in preparation for Run~3, the NSW is built with Micromegas (MM) and Thin Gap Chambers (sTGC), and supersedes the former Small Wheel instrumentation of MDTs, Cathode Strip Chambers (CSCs), and a single station of TGCs. The muon trigger system is provided, over the range $|\eta| < 2.4$, by Resistive Plate Chambers (RPCs) in the barrel ($|\eta| < 1.0$), while Thin Gap Chambers (TGCs) cover the endcap regions ($|\eta| > 1.0$), outside the toroid magnets. An additional contribution comes from the Micromegas (MM) and small-strip TGC detectors of the New Small Wheels, located between the cryostats of the endcap calorimeters and the endcap toroids. Interesting events are selected by a level-1 trigger system built in custom hardware, and are then filtered further by software algorithms running in the high-level trigger~\cite{ATLAScollab3}. Out of the 40~MHz bunch-crossing rate, the level-1 system retains events at a rate of at most 100~kHz; the high-level trigger applies additional selections so that data are written to disk at approximately 3~kHz. The configuration of the ATLAS detector during Run~3 is shown in \Cref{fig:ATLASexp}.
\\[2cm]
\begin{figure}[H]
    \centering
    \includegraphics[width=1\textwidth]{TESI/ATLASexp.png}
    \caption{Section of the Run 3 configuration of the ATLAS detector showing the locations of the larger detector sub-systems.}
    \label{fig:ATLASexp}
\end{figure}

\section{The Trigger System}
The ATLAS Trigger and Data Acquisition (TDAQ) system is responsible for selecting interesting events for analysis and reading them for offline processing, such as energetic leptons, photons, hadronic jets, tau leptons, or large missing energy. The TDAQ is organized into two selection levels: Level-1, made with custom electronics, and the High-Level Trigger, which is a software-based system running on commercial computers. These two levels share the Data Acquisition (DAQ) system, which carries data from the sub-detector electronics to offline processing following the trigger decisions. The L1 works on detailed information coming from the calorimeters and the muon spectrometer. Its maximum rate of acceptance sustainable by the readout systems is 100 kHz, while all processing has to be completed within the latency allowed by the detector electronics, which is \(2.5\ \mu s\) per event. The data acquisition system has also been upgraded so it can read the whole detector during Run 3. The ReadOut System needs to handle 100 kHz of events accepted by Level-1, plus an extra \(20\%\) margin. In Run 3, the load it needs to manage increases for two reasons: HLT asks for data more often, and the average event size goes up by \(30\%\), reaching \(2.1\ MB\) at \(\langle\mu\rangle \approx 60\). Further down, the system has to write an average flow of 3 kHz to mass storage while staying flexible enough to handle rate spikes. The max available bandwidth is 8 GB/s: in the example given, with \(\langle\mu\rangle \approx 60\) and a max physics event rate on disk of \(3.3\ kHz\), it needs \(6.9\), almost double compared to Run 2. In the end, to record raw data in continuous data acquisition at 3 kHz, at Point 1 there is at least \(1.4\ PB\) of effective space available: a capacity that ensures 24 hours of autonomy in case of downtime of CERN's mass storage services and accounts for both the time needed to transfer and then delete the files, as well as the margin required for the file system to maintain expected performance.

\clearpage

% ==================================================================
%  CAPITOLO 4 - MUON SPECTROMETER
% ==================================================================
\chapter{The ATLAS Muon Spectrometer}
\lettrine[lines=3,loversize=0.1,lraise=0.05]{H}{} igh momentum muons in the final state are among the most promising and robust signatures of the physics accessible at the Large Hadron Collider (LHC). Being the only charged particles that traverse the calorimeters almost undisturbed, they can be identified with little ambiguity even in the dense environment of \(pp\) collisions. To exploit this potential, the ATLAS Collaboration has designed a high-resolution Muon Spectrometer, able to trigger on muons and to measure their momentum in stand-alone mode over a wide range of transverse momentum, pseudorapidity and azimuthal angle. This chapter gives a brief description of the main components of the spectrometer, of the role each of them plays and of their arrangement within the system. Five detector technologies are employed for triggering and precision tracking: Resistive Plate Chambers, Thin Gap Chambers, Monitored Drift Tubes, Micromegas and small-strip TGCs, and each is presented in the following sections.


\section{Muon System}
The Muon Spectrometer (MS)~\cite{ATLASdesc} constitutes the outermost and largest part of the ATLAS experiment. Its role is to detect the charged particles that emerge from the barrel and endcap calorimeters, allowing us to measure their momentum up to \(|\eta| = 2.7\). The measurement relies on the deflection of muon tracks inside large superconducting air-core toroid magnets, complemented with two sets of chambers, one dedicated to triggering and one to high-precision tracking. Magnetic bending is provided by the large barrel toroid for \(|\eta| < 1.4\) and by two smaller endcap toroids, inserted at either end of the barrel magnet, for \(1.6 < |\eta| < 2.7\); in the transition region between them, muons are deflected by the combined barrel and endcap fields. This arrangement of magnets keeps the field orthogonal to the muon tracks, ensuring a smaller loss of resolution induced by multiple scattering. The MS joins the trigger system up to \(|\eta| = 2.4\). The MS is organised into a barrel and two endcaps. Three concentric stations of roughly cylindrical shape make up the barrel and are labelled: inner (BI), middle (BM) and outer (BO). Each endcap is, instead, built from three disc-shaped stations, known as wheels and labelled inner (EI), middle (EM) and outer (EO), to which an Extended Endcap ring (EE) is added at a radius larger than that of the endcap toroid cryostat. The three wheels occupy rather different positions along the beam line. The EI wheels (New Small Wheels) are housed within the volume enclosed by the barrel toroid, in the gap that separates the calorimeters from the endcap toroid cryostats. The EM wheels (Big Wheels) lie immediately downstream of those cryostats, while the EO wheels, with the same size of the EM ones, are carried by metal support structures anchored to the end-walls of the ATLAS cavern. The role of the EE rings is to add a third measurement point at intermediate pseudorapidity, \(1.05 < |\eta| < 1.3\), for the tracks that escape the coverage of the EO wheel; the rings are accordingly placed between the EI and EM wheels. A special case is that of the EIL4 stations, which cover the interval \(1.0 < |\eta| < 1.25\): instead of being carried by the EI wheels, they occupy a fixed position among the coils of the barrel toroid. As far as the barrel instrumentation is concerned, the bending coordinate is measured in all three stations by multilayer Monitored Drift Tube (MDT) chambers; the BM and BO stations carry, in addition, Resistive Plate Chambers (RPCs), from which both the trigger decision and the coordinate along the azimuthal direction are obtained. Only MDT chambers are installed in the outer wheels, while the middle wheels combine MDTs, for precision tracking in the bending plane, with TGCs, used both for triggering and for the azimuthal measurement. The inner wheels are located between the calorimeters and the endcap toroid cryostats, within the volume of the barrel toroid. They have been entirely replaced by the NSWs, which occupy the same location and cover the same polar range, \(1.3 < |\eta| < 2.7\). The NSWs use two chamber technologies: small-strip TGC (sTGCs) and micro-mesh gaseous structure (Micromegas) chambers. Both are fast enough to contribute to the Level-1 decision, and both are designed for precision tracking along the bending direction as well as for improved resolution in the azimuthal coordinate. The resulting gain in trigger \(p_T\) resolution and in background rejection makes it possible to retain a low muon \(p_T\) threshold at a sustainable Level-1 rate, thereby preserving the acceptance for many physics processes of interest. The layout of the MS is shown in \cref{fig:MS}, and its main parameters are summarised in \cref{tab:MS}.

\begin{figure}[H]
    \centering
    \includegraphics[width=1\textwidth]{TESI/MS_nuovo.png}
    \caption{Section of the ATLAS muon system, built around an air-core toroid magnet system, which measures the momentum of the charged particles leaving the barrel and endcap calorimeters within \(|\eta| < 2.7\); the components upgraded for Run 3, in particular the NSWs, are shown.}
    \label{fig:MS}
\end{figure}

\begin{table}[H]
    \centering
    \includegraphics[width=0.8\textwidth]{TESI/tab_04.pdf}
    \caption{Main parameters of the muon spectrometer}
    \label{tab:MS}
\end{table}

\section{Monitored Drift Tube (MDT)}
Monitored Drift Tube (MDT) chambers constitute a fundamental part of the precision tracking system of the ATLAS muon spectrometer, where they measure the trajectories of the charged particles produced in \(pp\) collisions over the pseudorapidity range \(|\eta| < 2.7\). They provide an average spatial resolution of about \(80\ \mu m\) for a single tube, which improves to roughly \(35\ \mu m\) once the measurements of a full chamber are combined. The chambers sustain hit rates as high as 300 kHz per tube with no appreciable loss of either spatial resolution or detection efficiency \cite{ATLASdesc, MDTTDCsim}.  Each chamber is equipped with up to 18 front-end boards, known as mezzanine cards. A mezzanine hosts three custom 8-channel Amplifier/Shaper/Discriminator (ASD) chips and one 24-channel Time-to-Digital Converter (TDC), the ATLAS Muon TDC (AMT), currently in its third version (AMT-3). The signals of all mezzanines belonging to a chamber are collected by an on-chamber data acquisition board, the Chamber Service Module (CSM), which formats and serialises them and transmits them over an optical link to the off-detector MDT Readout Driver (MROD) \cite{MDTTDCsim}. The read-out system was originally dimensioned for an instantaneous luminosity of \(10^{34}\ cm^{-2}s^{-1}\), while Run 3 has already exceeded \(2 \times 10^{34}\ cm^{-2}s^{-1}\) at a Level-1 accept rate of 100 kHz; dedicated simulations indicate that the legacy AMT-3 can cope with these conditions, the hit loss in the busiest chamber remaining below \(5\ \%\) \cite{MDTTDCsim}. For the HL-LHC, where the luminosity may reach \(7 \times 10^{34}~cm^{-2}s^{-1}\) and the Level-0 system is required to accept events at 1~MHz within a \(10\ \mu s\) latency, the front-end electronics will be replaced during Long Shutdown~3 and the new TDCs will operate without trigger mode.

\section{Resistive Plate Chamber (RPC)}
Muon identification in the barrel of ATLAS relies on Resistive Plate Chambers (RPCs) \cite{ATLASRPCRun3}, which deliver both a stand-alone trigger decision and a measurement of two coordinates. Each chamber consists of two independently operating detector layers, and the chambers are distributed over three concentric cylindrical shells located between roughly 7 and 10 m from the beam axis, immersed in a toroidal field of about 0.5 T. Every gas volume is built from bakelite electrodes enclosing a 2 mm gap, continuously flushed with a ternary mixture of \(C_2H_2F_4\), \(\text{iso - }C_4H_{10}\) and \(SF_6\) in proportions of \(94.7\%\), \(5.0\%\) and \(0.3\%\). Operation takes place in avalanche mode: the applied voltage of 9.6 kV corresponds to a field of 4.8 kV/mm across the gap, and it is rescaled automatically to compensate for deviations of temperature and pressure from the reference values of \(24\ ^\circ C\) and 970 mbar \cite{GA}. Signals are collected by readout panels whose strips, between 2.3 and 3.5 cm wide, are oriented so as to provide the two orthogonal coordinates \(\eta\) and \(\phi\). Altogether the system comprises 3720 gas volumes, spans 3650 \(m^2\) of active surface and is read out through approximately 385000 channels. A first upgrade of this layout has already been carried out: 16 RPC triplets (BIS78) were mounted in the transition region between barrel and end-cap, where the limited instrumentation in the interval \(1.0 < \eta < 1.3\) was responsible for a sizeable rate of fake triggers. In these new units, the bakelite electrodes are separated by a narrower 1 mm gap, and the copper readout strips are coupled to low-noise front-end electronics with high amplification. Further work is foreseen for the HL-LHC phase. The innermost barrel layer, at present equipped with drift tubes alone, will be complemented by a triplet of RPC layers, thereby extending the acceptance and adding redundancy to the muon trigger.

\section{Thin Gap Chamber (TGC)}
As discussed in section 4.1, the TGCs are located in the middle wheels, providing us with Level-1 trigger information and measuring the azimuthal coordinate of muons in the endcaps of the ATLAS MS \cite{TGC1}. TGCs are multi-wire proportional chambers with the characteristic that the wire-to-cathode distance  (1.4 mm) is smaller than the wire-to-wire distance (1.8 mm). The cell structure allows us to measure in a quasi-saturated mode due to a highly quenching gas mixture of \(CO_2\) and \(\text{n -} C_5H_{12}\) (n-pentane), i.e. with a gas gain of \(3\ \cdot\ 10^5\) that does not allow to make full use of its independence from the primary ionisation. Moreover, the highly quenching gas prevents the occurrence of streamers in all operating conditions. Instead, the presence of a high electric field around the TGC wires and, thanks also to the small wire-to-wire distance, leads to an incredible time resolution for the majority of the tracks. Including the variation of the propagation time on wires and strips, signals arrive with \(99\%\) probability inside a time window of 25 ns. More information on the fundamental parameters of TGC construction is shown in the TGC parameter book \cite{TGC}.

\section{New Small Wheel}
The NSW detectors allow us to reconstruct muon tracks with high precision and provide information for the L1 trigger \cite{ATLASdesc, NSW}. As discussed in Section 4.1, they use two peculiar chambers: small-strip TGC (sTGC) and micro-mesh gaseous structure (Micromegas) detectors. Both can be used for the trigger given the high ability of the sTGCs to identify single Bunch Crossings and the high speed of the MM thanks to their small drift gap. Also, they are great for tracking because of the high resolution (\(100\ \mu m\) in the bending direction. These two different detectors reach high spatial resolution with different mechanisms: the MM exploits a very thin strip pitch (\(<0.5\ mm\)), while the sTGCs, having a larger pitch (\(3.2\ mm\)), rely on the charge sharing between the strips. This precision needs to preserve the high muon momentum resolution even in a highly background enviroment as LHC. Moreover, improve a better resolution for the azimuthal coordinate with respect to the EI wheels: better resolution at trigger level and offline reconstruction, in view of Run 3. All this leads to measuring the transverse muon momentum in the endcaps that, combining the track segments in the NSW with those from the ID, must have a better precision of \(15\%\) for a 1 TeV muon all over \(|\eta|<2.7\). To do this, NSW must have a resolution of \(50\ \mu m\) in the bending plane to reconstruct these track segments, independently of the background conditions. Since the NSW has 16 sectors, this leads to a resolution of \(150\text{-}175\ \mu m\) for each strip plane. In fact, from a geometric point of view, each wheel maintains the spectrometer's eightfold symmetry and is divided into eight Large sectors and eight Small sectors (see \cref{fig:NSWa}, instead, the \cref{fig:NSWb} shows the shielding disk JD). In the end, the layout reflects the division of tasks between the two technologies: the sTGCs are on the outside because that's where their angular precision is most useful for the trigger, while the Micromegas occupy the central volume because they're responsible for precise tracking measurement in the curvature coordinate.

\begin{figure}[htbp]
    \centering
    \begin{subfigure}[b]{0.48\textwidth}
        \centering
        \includegraphics[width=\textwidth]{TESI/NSWa.png}
        \caption{}
        \label{fig:NSWa}
    \end{subfigure}
    \hfill
    \begin{subfigure}[b]{0.48\textwidth}
        \centering
        \includegraphics[width=\textwidth]{TESI/NSWjd.png}
        \caption{}
        \label{fig:NSWb}
    \end{subfigure}
    \caption{New Small Wheels: \textit{(a)} shows the structure: the wheel on the left is seen from the cavern end-wall, with the A-frame and shielding hub visible, along with the eight Large sectors and the eight spokes supporting them; the wheel on the right is seen from the IP, and since the newJD shielding disc is not shown, the eight Small sectors, aligned with the barrel toroid coils, are visible, along with the spokes supporting them. \textit{(b)} Support structures for the NSWs, including the newJD discs. }
    \label{fig:nsw-layout}
\end{figure}

\subsection{small-strip TGC (sTGC)}
The sTGCs \cite{ATLASdesc, NSW} were created to do two things at once: provide tracking fast enough to enter the level 1 hardware trigger and, at the same time, a position measurement precise enough to also be used in offline muon reconstruction. Classic TGCs were only used for the trigger; the new thing here is that the same detector has to do both jobs, in an environment with higher background. The technological base remains that of the TGCs: a \(2.8\ mm\) gap bordered by two FR4 cathodes coated with a resistive mixture of graphite and resin, filled with a mixture of \(55\%\ CO_2\) and \(45\%\) n-pentane; in the center of the gap, the anode plane consists of \(50\ \mu m\) gold-plated tungsten wires spaced \(1.8\ mm\) apart. The changes, on the other hand, only concern what was needed for the two new goals. The wires are now parallel to the radial axis of the chamber, whereas in the TGC they are stretched in the azimuthal direction; the readout strips are much denser (\(2.7\ mm\) wide, \(3.2\ mm\) pitch), with a geometry experimentally optimized to limit capacitive coupling between channels and nonlinear corrections, while keeping the number of readout channels as low as possible; finally, the signal layer has been moved under a thin pre-preg, just below the graphite instead of on the back of the FR4, increasing the capacitance to maintain the same transparency to fast signals. All of these choices result in an intrinsic resolution of about \(45\ \mu m\) for particles hitting perpendicularly \cite{NSW2}. To deal with background, they adjusted the resistivity of the cathode coating, lowering it to \(150\ k\Omega/\Box\)\footnote{Not to be confused with resistance, surface resistivity is the resistance between two opposite sides of a square of a thin-film resistor, and it doesn't depend on the size of the square; it's, anyway, measured in \(\Omega\).} in the chambers closest to the interaction point and to \(200\ k\Omega/\Box\) in the others, so the accumulated charge is quickly dispersed. Together with reducing the decoupling resistors, this keeps the chambers efficient for minimally ionizing particles even at rates up to \(20\ kHz/cm^2\) over large areas.

\subsection{micro-mesh gaseous structure (Micromegas)}

Micromegas \cite{ATLASdesc} are thin planar gas detectors in which the traditional high-voltage wire planes are replaced by a metal micro-mesh. The detector consists of a drift electrode, a gas gap of a few millimetres acting as conversion and drift region, and the micro-mesh, placed \(120\text{-}130\ \mu\text{m}\) above the readout strips, which defines the amplification region. The drift field (a few hundred V/cm) is set close to the plateau of the drift velocity to minimise longitudinal diffusion and ion backflow, while the amplification field (\(40\text{-}50\ \text{kV/cm}\)) provides a gas gain of about \(10^4\). Since the amplification field is 50--100 times stronger than the drift field, the mesh is transparent to more than \(95\%\) of the drifting electrons. Electron drift is relatively slow (up to about \(100\ \text{ns}\) in ATLAS), whereas the avalanche develops in a fraction of a nanosecond, producing a very fast signal on the strip. The ions, mostly created in the last multiplication stage close to the strip, are quickly collected by the mesh. The main challenge for large-area detectors is discharges, which occur when an avalanche exceeds a few tens of millions of electrons (Raether limit) and can damage both the detector and the readout electronics. Given the gain of about \(10^4\) required for minimum ionising muons, any ionisation releasing more than a thousand electrons within a region comparable to the lateral size of the avalanche can trigger them. To mitigate this, a layer of resistive strips is placed above the readout strips, separated by a thin insulator, so that the readout is capacitively coupled to the amplification region. This costs a fraction of the signal amplitude, but allows operation at higher gain and reduces the discharge intensity by about three orders of magnitude.



%===================================================================
% CAPITOLO 5 - RICOSTRUZIONE E CLASSIFICAZIONE MUONI
%===================================================================
\chapter{Muon Reconstruction \& Classification}


\lettrine[lines=3,loversize=0.1,lraise=0.05]{I}{n} ATLAS, muon reconstruction is based on information from several detector systems, including the Inner Detector (ID), the Muon Spectrometer (MS), and the calorimeters~\cite{Journ}. Muons are highly penetrating particles (i.e. minimum-ionising particles), identified using tracks in the MS. To ensure high efficiency and precision, the muon reconstruction combines the information of the different subdetectors. This chapter describes how the global reconstruction works, starting from the stand-alone tracks and from the different muon types that derive from them (Sec. 5.2). The data samples used in the analysis are then described (Sec. 5.3). Finally, the chapter presents the first part of the analysis in Sec 5.4: the selection of prompt muons only and the application of the working points (Sec.~5.4) to study how the composition by muon type changes. The goal is to obtain a sample of CB muons only, on which the resolution measurement of the following chapters is based.

\section{Stand-alone tracks reconstruction}
The stand-alone track reconstruction in the MS is a fundamental process to obtain the characteristics and trajectories about the muon that pass through the detector. This process occurs in several steps that involve combinations of geometrical and physical methods to guarantee a high quality and reliable muon track reconstruction. It starts with the short straight-line local track segments derived from "hits" in an individual MS station (there are typically 3 MS stations in most of the spectrometer). Each segment represents a part of the muon's trajectory, and their identifications are performed using the Hough transform. The trigger chambers provide the second coordinate (along the non-bending direction) to the segments reconstructed in each station, so that segments from different stations can be matched in both coordinates. Once these segments are identified in the different MS stations, they are combined to form preliminary candidate muon tracks using a loose pointing constraint based on the interaction point (IP) and the parabolic trajectory that constitutes a first-order approximation to the muon bending in the magnetic field. A global \(\chi^{2}\) fit is then performed on each track candidate, taking into account the magnetic field. This statistical technique minimises the residuals between the measured hits in the MS and the expected trajectory, which can be affected by misalignments between the detector chambers as well as by the interactions of the muon with the detector material, such as multiple scattering and energy loss, that alter its real path. From the fit, two operations start, that is the removal of outlier hits and the addition of the hits along the trajectory that were not assigned to the original track candidate in order to improve the reconstruction. Last but not least, if some ambiguities are present between tracks, for example, two tracks that share the same hits fraction, the lower-quality one is rejected. One exception to this rule is made in the case of tracks that are identical in two stations but share no hits in a third station. This helps to guarantee a high efficiency in the light-flavour muons system reconstruction tracks originating from a second vertex due to the decay of pions or kaons, where the detector can mistakenly interpret it as a true signal from a prompt muon track. The final set of tracks is re-fitted with a loose IP constraint, taking into account the energy loss in the calorimeters, and back-extrapolated to the beam line. The \(p_{T}\) of the extrapolated track is then expressed at the IP.

\section{Global muon reconstruction}
The global muon reconstruction is based on a system that involves all the detectors in ATLAS experiment, as the ID, the MS and calorimeter systems. Also here, one wants to achive the same result, i.e. trajectory description and properties, but improved by all the info in the detector. There are five strategies to obtain the reconstruction, each producing different muon types. Combined Muons (CB) are reconstructed by comparing tracks that come from MS with those coming from ID. Then a fit is made using the MS and ID hits, taking into account the energy loss in the calorimeters, which can influence the muon path. For \text{$|\eta| > 2.5$}, muon tracks in MS can be associated with partial information from Pixel and SCT detectors, creating a subcategory of CB muons called Forward Muons associated with Silicon Forward muons (SiAF) that extend the coverage beyond the ID acceptance. Inside-Out muons (IO), where the tracks are reconstructed starting from ID and then extrapolated to MS, improving the efficiency in regions where MS is limited and for muons that have low \( p_T \). If an MS track cannot be matched to an ID track, its parameters are extrapolated to the beamline and used to define an Extrapolated Muon (ME), that are just the Stand Alone muons (SA). This expands the coverage of the MS to \text{$|\eta| = 2.7$}, improving the muon identification in regions where ID failed. Segment-Tagged Muons (ST) are reconstructed starting from an ID track, which is identified as a muon because it is associated with at least one segment in the MS. Unlike the previously mentioned categories, ID and MS information are not combined here; only ID information is used, while MS is used as support to select the tracks. ST muons are used to recover some efficiency in regions where the MS has limited acceptance, and it's not possible to reconstruct even a partial track, such as the gap at $|\eta|=0$ left by the passage of the detectro services. Calorimeter-Tagged Muons (CT), as the name suggests, are muons reconstructed starting from ID tracks, which are associated with energy loss in the calorimeters compatible with minimum ionising particles. Also in this case, only the Inner Detector (ID) is used to reconstruct the track, while the calorimeter serves as support to recover efficiency in regions where the Muon Spectrometer (MS) has no acceptance.
\clearpage
\section{Data Samples}

The analysis presented in this chapter is based on Monte Carlo simulated samples: the studied observable requires a comparison between reconstructed momentum and the truth momentum of the muons. All the simulated processed refers to \(pp\) collisions at \(\sqrt{s}=13.6\ TeV\), generated to reproduce the data-taking detector conditions of the ATLAS detector during the LHC Run 3 with the pile-up conditions relating to 2024. The response of the detector is simulated with the toolkit \textsc{GEANT4} \cite{geant4} and the events are reconstructed with the same software chain applied to the real data. Two physics process categories are implemented in this analysis. The first one is the Drell-Yan production of the Z bosons with decay into two muons, $Z\rightarrow\mu^+\mu^-$ (DSID 601190), generated with \textsc{Powheg Box} at the Next-to-Leading-Order (NLO) accuracy in QCD and \textsc{Pythia 8} for matrix element (ME) calculations, for the modelling of the parton shower, hadronisation, and the underlying event, with the parameter sets AZNLO; the effect of QED final-state radiation was simulated with \textsc{Photos++}. This sample provides us with high isolated muon statistics with high quality in the region of the Z peak (i.e. \(p_{\mathrm{T}}\simeq45\ GeV\)), that is, the region in which the ATLAS collaboration determines the muon calibration constants \cite{ATLAS-MCP-2023}. The second one is made up of a heavy resonance of \(Z'\rightarrow\mu^+\mu^-\) at zero width (DSID 801862-801866), generated with \textsc{Pythia 8} using a parton density function (PDF) set at the Leadin-Order (LO) level of accuracy and the A14 set of tuned generator parameters were used for five mass values: 500, 1000, 3000, 5000 and 8000 GeV as shown in \cref{fig:peaks} and \cref{tab:peaks}. 
\begin{figure}[htpb]
    \centering
    \includegraphics[width=0.85\textwidth]{TESI/plot_mmumu_Z_Zprime_pieni_minimal.pdf}
    \caption{Distribution of the invariant mass $M_{\mu\mu}$ of the muon pair in the
  $Z\rightarrow\mu^{+}\mu^{-}$ sample and in the five $Z'$ samples with different masses. Each histogram is normalized to its own integral and every distributions has peaks at the mass of the corresponding resonance. The small excess in the $Z$ sample at $M_{\mu\mu} \approx 10$~GeV is due to pairs formed by a muon from the $Z$ decay and a soft non-prompt muon from the rest of the event, when the second muon from the $Z$ is not reconstructed.}
    \label{fig:peaks}
\end{figure}

\begin{table}[htpb]
    \centering
    \begin{tabularx}{\textwidth}{X|X|X|X|X}
        \toprule
        \textbf{Sample} &
        \textbf{Resonance [GeV]} &
        \textbf{Peak $M_{\mu\mu}$ [GeV]} \\
        \midrule
        $Z \rightarrow \mu^+\mu^-$ & 91.2   & $94 \pm 5$ \\
        $Z'$ (500 GeV)             & 500    & $497 \pm 26$ \\
        $Z'$ (1000 GeV)            & 1000   & $1030 \pm 54$ \\
        $Z'$ (3000 GeV)            & 3000   & $2920 \pm 150$ \\
        $Z'$ (5000 GeV)            & 5000   & $4910 \pm 260$ \\
        $Z'$ (8000 GeV)            & 8000   & $8270 \pm 430$ \\
        \bottomrule
    \end{tabularx}
    \caption{Mass of each simulated resonance and position of the peak of the corresponding
  $M_{\mu\mu}$ distribution shown in \cref{fig:peaks}, taken as the centre of the bin and the uncertainty half the bin width. In all samples the observed peak is compatible with the resonance mass.}
    \label{tab:peaks}
\end{table}
\noindent Since the resonance width is zero, each sample provides us with two muons with fixed momentum to \(m_{Z'}/2\) in the resonance reference frame at rest; in the lab reference frame, the transverse muon momentum distribution presents the usually jacobian peak at that value, allowing us to have a full popolate spectrum since the \(Z\rightarrow\mu^+\mu^-\) process have not enough statistics\cite{ATLAS-MCP-2023}, as shown in \cref{fig:mass_statistics}. 
\begin{figure}[htpb]
    \centering
    \includegraphics[width=0.85\textwidth]{TESI/plot_pt_Z_Zprime_completo_conteggi_minimal.pdf}
    \caption{Number of reconstructed muons per bin as a function of the muon $p_{\mathrm{T}}$. The simulated $Z\rightarrow\mu^{+}\mu^{-}$ sample provides the statistics up to about 200 GeV; above this value, the simulated $Z'\rightarrow\mu^{+}\mu^{-}$ samples at different masses are needed to populate the high-$p_{\mathrm{T}}$ region, from 200~GeV up to the TeV scale.}
    \label{fig:mass_statistics}
\end{figure}
Each process was built in two different configurations of the detector geometry. The first, in the following denoted as \emph{ideal} geometry, assumes a perfect alignment of both the ID and the chambers of the MS: this is the standard condition for the ATLAS simulations, in which no residual misalignment effect appears. The second, denoted as \texttt{MSmisalign}, the chamber positions in the intermediate layer of the MS were deliberately shifted by an uniform amounts compatible with the residual precision achieved by the alignment procedures \cite{ATLAS-MCP-2023}, reproducing a realistic chamber shift, which is a distorted geometry that shows their misalignment. The comparison between the two configurations, given the same generated process and applied selection, is the method used in Section 7.3 to isolate the contribution of the residual misalignment. From a technical view, the simulated data samples used in this analysis are \emph{ntuples} \texttt{ANALYSIS.root} (production with \texttt{MCP\_TESTNTUP}), given by files AOD through the package \texttt{MuonxAODAnalysis}. These ntuples contain, for each selected muon, kinematic characteristics reconstructed separately from the track in the ID, in the MS, and from the combined track (CB). The $Z\rightarrow\mu^+\mu^-$ sample, consisting of 68 files with a perfectly aligned MS chamber configuration, provides sufficient statistics up to \(\sim 200\) GeV, while the $Z'$ samples, consisting of 26 files with a misaligned MS chamber configuration, extend the coverage up to the TeV scale, for a total of \(\simeq 400\) million events.

\section{MS chamber alignment}
The alignment of the MS chambers is crucial for a good muon reconstruction, since any displacement of the chambers with respect to their nominal positions directly translates into a bias on the measured momentum. The goal is to know the position of the chambers with a precision of about \(50\ \mu\text{m}\) over distances of the order of 10 m, while the actual displacements can reach a few millimetres; for instance, shifts up to 2.7 mm with respect to the nominal position have been observed in the NSW chambers when the toroid is switched on. The alignment is performed in two complementary ways. An optical alignment system, based on RASNIK sensors (a CCD camera imaging a coded mask through a lens), measures the relative displacements of the chambers between different alignment runs. This is complemented by a track-based alignment, which uses special runs with the magnetic field off, where muon tracks are straight lines, and exploits the ID to select high-\(p_T\) tracks in order to reduce the effect of multiple scattering. Some deformations, the so-called weak modes, such as common shifts of the chambers and radial distortions, leave the tracks compatible with a straight line and are therefore only partially recovered, using the overlaps between adjacent sectors. In Run 3, the alignment reaches a precision on the sagitta of about \(50\ \mu\text{m}\) in the barrel and about \(100\ \mu\text{m}\) in the endcaps \cite{ATLAS_MS_alignment_2021, Rasnik_ultimate}.


\chapter{Muon momentum resolution}

\lettrine[lines=3,loversize=0.1,lraise=0.05]{T}{} he precision with which muon momentum is reconstructed directly determines the resolution of the invariant mass of every final state that contains muons, and thus the sensitivity of precision measurements and of searches for new physics in these channels. In this chapter, the momentum resolution is measured on a simulated sample through the curvature relation discussed in Section 7.2 in a transverse momentum range going from the region of the \(Z\) peak to the TeV scale. Since the \(Z\to\mu^+\mu^-\) process does not have enough statistics at high \(p_{\mathrm{T}}\), an additional \(Z'\) sample is used to populate the whole spectrum. The measurament is repeated for an ideal geometry configuration where all the MS chambers are in a perfectly aligned, and for a configuration where the chambers are completely misaligned by a uniform quantity. To do this, a parametrization \cite{ATLAS-MUON-2016, ATLAS-MCP-2023}, discussed in the next sections, which expresses the resolution as the sum in quadrature of three terms with different dependence on the transverse momentum of the muon. Since the terms that increase linearly with \(p_T\) is sensitive to residual misalignments between the elements of the ATLAS detector, comparing samples produced with ideal geometry and with the spectrometer deliberately misaligned allows us to isolate and quantify the contribution of residual misalignment to the degradation of high-\(p_T\) resolution. The analysis is developed on \texttt{lxplus} \cite{lxplus}, a CERN interactive login service.



\section{Prompt muon selection}

Muons can be divided into two classes: \emph{Non-prompt muons} which are classified into light flavour muons, coming from the decay of light hadrons such as pions and kaons produced in a gluon jet, and heavy flavour muons, coming from the decay of B- and C-hadrons; \emph{Prompt muons} which are produced directly at the primary vertex and are more isolated in the ID and/or the calorimeters, leave a clean track that can be easily reconstructed and it comes principally from the \(W\), \(Z\) and \(Higgs\) boson decays, or from other electroweak processes. In ATLAS Truth Classification, which is a tool that combines the type and origin truth of the particle to separate these kind of  muons produced at the primary vertex from those generated in hadron decay or in the interactions with the materials of the detector, \footnote{\url{https://gitlab.cern.ch/atlas/athena/-/blob/main/PhysicsAnalysis/AnalysisCommon/TruthClassification/README.md\#prompt-muons}} prompt muons are those that satisfy \(\text{truthmuon\_IFFType} == 4\), which correspond to the category PromptMuon: it can be only if his \texttt{truthType} is that of an isolated muon and its \texttt{truthOrigin} is a prompt source. Without this constraint the selected sample could be contain also muons coming from heavy hadron decay, produced at distances of the order of centimeters from the primary vertex and classified separately as \(\text{IFFType} == 8\) for B-hadron and \(\text{IFFType} == 9\)  for C-hadron; in-fly-decay muons coming from pions and kaons within the detector, that give rise to a \emph{kink} in the track and they fall into the category of light flavour, \(\text{IFFType} == 10\); muons which truth type, instead, correspond to an electron, \(\text{IFFType} == 12\); and in the end muons which they're tool can't led to any category, or due to missing truth information \(\text{IFFType} == 1\), or because can't be associated to anyone class we juste define, \(\text{IFFType} == 0\), called \texttt{Unkonwn}, which also belong to this category fake muons produced from bad reconstructed tracks.

\section{Curvature resolution applied on Prompt Muon}

We begin with the simulated $Z\rightarrow\mu\mu$ sample alone, in the ideal configuration. The quantity on which the first part of this analysis is based is the curvature resolution $\mathcal{R}$ defined as the relative difference between the reconstructed curvature and the truth one:

\begin{equation}
  \mathcal{R} \;\equiv\;
  \frac{1/p_{\mathrm{T}}^{\mathrm{reco}} - 1/p_{\mathrm{T}}^{\mathrm{truth}}}
       {1/p_{\mathrm{T}}^{\mathrm{truth}}}
  \;=\;
  \frac{p_{\mathrm{T}}^{\mathrm{truth}}}{p_{\mathrm{T}}^{\mathrm{reco}}} - 1 .
  \label{eq:osservabile}
\end{equation}

\noindent and will be positive when the reconstructed momentum underestimates the truth. We work with curvature resolution and not directly with the momentum because \(1/p_{T}\) is proportional to the sagitta, that is the quantity measured by the detectors \cite{Gluckstern1963}. For each reconstructed muon, we can obtain information about the truth muon by the associated index in the ntupla. One requirement is asked about this muons, that this they have to be \emph{prompt} according to the truth classification \texttt{(truthmuon\_IFFType == 4}). No requirements will be applied to the reconstructed muon. The analysis will be done by comparing both prompt and non-prompt muons, in order to show the effects. The \(R\) distribution is filled in 100 intervals between \(-0.2\) and \(+0.2\), in a inclusive form and in 7 intervals of \(p_T^{truth}\): $0\text{-}20$, $20\text{-}24$, $24\text{-}27$, $20\text{-}24$, $27\text{-}30$, $30\text{-}33$, $33\text{-}37$, $37\text{-}40$, $40\text{-}42$, $42\text{-}44$, $44\text{-}46$, $46\text{-}50$, $50\text{-}60$, $60\text{-}90$ and $90\text{-}200$ GeV. Intervals are widen as the momentum increases to compensate for the low statistics of the sample \(Z\rightarrow\mu^+\mu^-\), passing from tens of millions of muons in the first intervals to hundreds thousand to the latter intervals. Prompt muons make up \(99.25\%\) of the sample\(Z\) in the ideal geometry, leading to a small difference in RMS (\(\simeq1\%\), almost invisible to the naked eye, \cref{fig:res-inclusiva}). 
\begin{figure}[htpb]
  \centering
  \includegraphics[width=\textwidth]{TESI/h_res_prompt_vs_incl_minimal.pdf}
    \caption{Distribution of the curvature resolution $R$ of \cref{eq:osservabile} for muons of the $Z \to \mu^+\mu^-$ sample in the ideal geometry, comparing the inclusive selection with the one restricted to prompt muons (\texttt{truthmuon\_IFFType == 4}). Prompt muons make up $99.25\%$ of the sample, so the two distributions are visually almost indistinguishable and their RMS differ by about $1\%$, $(28.778\pm0.001)\cdot10^{-3}$ against $(28.445\pm0.001)\cdot10^{-3}$. The effect of the non-prompt contamination is not uniform in $p_{\mathrm{T}}$: as shown in \cref{tab:nonprompt_fraction}, it is concentrated in the lowest bin, where non-prompt muons reach $5.3\%$ of the sample and inflate the RMS, while above 20~GeV it drops below $0.1\%$ and the two values coincide to the third decimal digit.}
  \label{fig:res-inclusiva}
\end{figure}
\noindent However, looking at how that \(0.75\%\) of non-prompt muons is distributed per \(p_T^{truth}\) bin, the story is different as shown in \cref{tab:nonprompt_fraction}. Non-prompt contamination is almost focused on the fisrt bin (0-20 GeV, \(5.3\%\)) then start become under \(0.1\%\) and both RMS are egual to the third decimal digit. This is consistent with the fact that prompt muon coming from \(Z\rightarrow\mu^+\mu-\) has a peak in the spectrum around \(m_Z/2\approx45\ GeV\), while non-prompt muon comes from \(\pi/K\) in-flight decay (where the decay happens inside the detector and the trajectory ends up being deflected) or heavy hadron decay has softer spectrum below 20-30 GeV \cite{ATLAS-MUON-EFF-2021}. In the 0-20 GeV isolated bin: the prompt RMS, $(27.289\pm 0.003) \times 10^{-3}$, is practically identical up to 37 GeV bins, showing that the apparent worsening at low pt is not due to the detector resolution but only due to contamination from non-prompts. The quantitative trend is summarized by the RMS of each $p_{\mathrm{T}}^\mathrm{truth}$ interval in \cref{fig:rms-vs-pt}, where we can already see three regimes of the equation \eqref{eq:parametrizzazione}, highly discussed in the next chapter. 
\begin{table}[htbp]
  \centering
  \begin{tabular*}{\textwidth}{@{\extracolsep{\fill}}cccc}
    \toprule
    $p_{\mathrm{T}}^{\mathrm{truth}}$ [GeV] & \% non-prompt & RMS inclusive & RMS prompt \\
    \midrule
    0--20    & 5.30\%   & $29.752 \pm 0.003$ & $27.289 \pm 0.003$ \\
    20--24   & 0.091\%  & $26.770 \pm 0.004$ & $26.761 \pm 0.004$ \\
    24--27   & 0.046\%  & $27.064 \pm 0.004$ & $27.060 \pm 0.004$ \\
    27--30   & 0.028\%  & $27.345 \pm 0.004$ & $27.344 \pm 0.004$ \\
    30--33   & 0.017\%  & $27.646 \pm 0.004$ & $27.645 \pm 0.004$ \\
    33--37   & 0.010\%  & $27.975 \pm 0.003$ & $27.974 \pm 0.003$ \\
    37--40   & 0.007\%  & $28.348 \pm 0.004$ & $28.347 \pm 0.004$ \\
    40--42   & 0.005\%  & $28.579 \pm 0.004$ & $28.578 \pm 0.004$ \\
    42--44   & 0.004\%  & $28.825 \pm 0.004$ & $28.825 \pm 0.004$ \\
    44--46   & 0.003\%  & $29.021 \pm 0.005$ & $29.021 \pm 0.005$ \\
    46--50   & 0.004\%  & $29.287 \pm 0.004$ & $29.287 \pm 0.004$ \\
    50--60   & 0.007\%  & $29.945 \pm 0.004$ & $29.944 \pm 0.004$ \\
    60--90   & 0.008\%  & $31.913 \pm 0.006$ & $31.912 \pm 0.006$ \\
    90--200  & 0.006\%  & $37.260 \pm 0.013$ & $37.260 \pm 0.013$ \\
    \midrule
    0--200   & 0.75\%   & $28.778 \pm 0.001$ & $28.445 \pm 0.001$ \\
    \bottomrule
  \end{tabular*}
  \caption{Fraction of non-prompt muons and RMS of the curvature resolution $R$, inclusive and for prompt muons only, in each $p_{\mathrm{T}}^{\mathrm{truth}}$ bin of the $Z \to \mu^+\mu^-$ sample in the ideal geometry. The uncertainties are obtained with the ROOT method \texttt{GetRMSError}, corresponding to $\sigma_{\mathrm{RMS}} = \mathrm{RMS}/\sqrt{2N}$. The non-prompt contamination is concentrated in the lowest bin, where it reaches $5.3\%$ and visibly inflates the inclusive RMS, while above 20~GeV it falls below $0.1\%$ and the two columns coincide within the quoted uncertainties. The prompt RMS in the first bin, $(27.289\pm 0.003) \times 10^{-3}$, differs by less than $1\%$ from the one of the 20-37~GeV bin, showing that the apparent degradation at low $p_{\mathrm{T}}$ is due to the non-prompt contamination and not to the detector resolution.}
  \label{tab:nonprompt_fraction}
\end{table}
\begin{figure}[htpb]
  \centering
  \includegraphics[width=\textwidth]{TESI/rms_vs_pt_minimal.pdf}
    \caption{RMS of the curvature resolution $R$ as a function of $p_{\mathrm{T}}^{\mathrm{truth}}$ for prompt muons of the $Z \to \mu^+\mu^-$ sample in the ideal geometry, with the values listed in \cref{tab:nonprompt_fraction}. Each point is placed at the mean value of $p_{\mathrm{T}}^{\mathrm{truth}}$ bin, and the horizontal bars indicate the bin width; the vertical bars are the statistical uncertainties $RMS/\sqrt{2N}$ and are too small to be visible. The trend already shows two of the three regimes of \cref{eq:parametrizzazione}, highly discussed in the next chapter.}
  \label{fig:rms-vs-pt}
\end{figure}
\clearpage

\subsection{Pseudorapidity $\eta$ dependence}

An extension of the analysis resolution in $p_{\mathrm{T}}$ regards how the resolution has a dependence on $|\eta^{\mathrm{truth}}|$. We analyze the same quantity discussed in Section 5.4.1 for prompt muons, but for both $p_{\mathrm{T}}^{\mathrm{truth}}$ and $|\eta^{\mathrm{truth}}|$ in 6 $|\eta^{\mathrm{truth}}|$ bins that follow the geometry of the ATLAS detector. For each bin, we can extrapolate the RMS, with uncerntanty evaluated as $RMS/\sqrt{2N}$, where N is the number of enteries. The result, shown in \cref{fig:map_eta}, indicates which region has the best resolution and coverage for working with clean tracks. The vertical axis reports the 14 $p_{\mathrm{T}}$ intervals used for a first analysis of the $Z\to\mu\mu$ sample in an ideal geometry of the MS chambers, while the horizontal axis reports the $\eta$ intervals, which show the coverage. The coloured scale, going from white to blue, shows an RMS that increases gradually towards the regions where the ID no longer has coverage ($|\eta|>2.5$) and as the muon momentum increases. 
\begin{figure}[htpb]
    \centering
    \includegraphics[width=\textwidth]{plot_RMS_eta_pt_map_minimal.pdf}
    \caption{RMS of the curvature resolution as a function of $|\eta^{\mathrm{truth}}|$ for prompt muons of the $Z \to \mu^+\mu^-$ sample in the ideal geometry, shown separately for the 14 $p_{\mathrm{T}}^{\mathrm{truth}}$ bins.} 
    \label{fig:map_eta}
\end{figure}
\noindent The resolution is best in the region \(0.1 \le |\eta^{\mathrm{truth}}| < 1.01\). In the innermost bin, \(|\eta^{\mathrm{truth}}| < 0.1\), the gap in the MS, reserved for the passage of services, reduces the coverage and the resolution gets worse at high \(p_{\mathrm{T}}^{\mathrm{truth}}\). The resolution also degrades in the transition region between the barrel and endcap toroids (\(|\eta^{\mathrm{truth}}| \in [1.3, 1.7]\)), due to the low and reversing magnetic field~\cite{ATLAS-MCP-2023}. The outermost bin, \(2.5 \le |\eta^{\mathrm{truth}}| < 2.7\), beyond the acceptance of the Inner Detector, is almost independent of \(p_{\mathrm{T}}^{\mathrm{truth}}\)~\cite{ATLAS-MUON-EFF-2021}.

\section{Working Point}

Quality cuts called Working Point (WP) can be applied. WPs aim to reject light muons, which generally produce bad-quality muons since, having been produced by a secondary decay, still in flight in the detector before becoming muons in the final state, they can cause a fake track that the detector confuses with a prompt muon track. Instead, muons coming from heavy flavour hadrons, like B-hadrons or C-hadrons, produce high-quality muons in the final state that can be easily distinguished from Prompt Muons. WPs are characterised by their purity and efficiency. The efficiency of a prompt muon selection is the probability that one muon that flies through the detector is reconstructed as a real muon that satisfies all the WP criteria. Similarly, the purity of a WP selection is given by one minus the misidentification hadron rate, where this one represents the light hadron fraction reconstructed as a real muon that satisfies the WP criteria. Three WPs are adopted in order to increase purity and decrease efficiency: Loose, Medium and Tight. The Loose WP was designed for the reconstruction of Higgs boson decays in the four-muon final state, which, due to the high muon multiplicity and large signal-to-background ratio, benefits from a higher efficiency at the cost of less purity and larger systematic uncertainties. Medium WP has a good purity-efficiency ratio, maintaining lower systematic uncertainties in both prompt muon identification efficiency and light flavour background cuts, and is suitable for a wide range of analyses. Tight WP selection has the highest purity in terms of light flavour background rejection and applies all the quality cuts used in the Medium WP, along with additional criteria. Identification criteria, selected muons with respect to the number of measurement stations, fit quality, and compatibility with independent ID and MS measurements \cite{ATLAS-MUON-EFF-2021}. Now it's right to ask whether the resolution depends on the muon WPs quality. In order to do this, were compared the three WPs discussed: \emph{Loose}, \emph{Medium} and \emph{Tight} (that actually are cumulative, i.e. Tight implies Medium that implies Loose). The three distribution as shown in \cref{fig:wp-overlay} are indistinguishable, meaning that the choice of a WPs does not influence the $p_{\mathrm{T}}$ resolution, but only the efficiency. 
\begin{figure}[htbp]
  \centering
  \includegraphics[width=0.75\textwidth]{TESI/h_res_wp_overlay_minimal.pdf}
    \caption{Normalised distributions of the curvature resolution $R$ of \cref{eq:osservabile} for prompt muons of the $Z \to \mu^+\mu^-$ sample in the ideal geometry, compared for the three identification criteria: Loose, Medium and Tight. The three working points are cumulative, Tight implying Medium and Medium implying Loose. The Loose and Medium curves are indistinguishable and give the same RMS, while Tight is smaller. The choice of working point therefore affects the selection efficiency, quantified in \cref{fig:efficiency} (a), (b) and (c), rather than the momentum resolution itself.}
  \label{fig:wp-overlay}
\end{figure}
Since the efficiency is here a WP selection efficiency, defined as the number of muons that satisfied the WP criteria with respect to the reconstructed muons that are already matched with a truth muon, we obtain $98.7\%$ for Loose, $97.7\%$ for Medium, and $90.4\%$ for Tight and errors evaluated with the Clopper-Pearson method (\emph{TEfficiency}), as shown in \cref{fig:efficiency}. Besides efficiency, an important aspect for evaluating the performance of identification working points is the purity of the sample being used, defined as the fraction of reconstructed muons that meet the working point criteria and are matched to a truth muon. While efficiency measures the loss of real muons, purity measures the contamination from objects not matched to truth muons, such as products from in-flight decays of pions and kaons. In \ref{fig:purezza_wp} are shown the purity as a function of $p_mathrm{T}$, $eta$, and $phi$ for the Loose, Medium, and Tight working points. Overall, the stricter criteria of the Tight working point ensure the highest purity, at the cost of lower efficiency.
\begin{figure}[htbp]
    \centering
    \includegraphics[width=\textwidth]{TESI/efficiency2.pdf}
    \caption{Identification efficiency of the Loose, Medium and Tight working points as a function of (a) $p_{\mathrm{T}}^{\mathrm{truth}}$, (b) $\eta^{\mathrm{truth}}$ and (c) $\varphi^{\mathrm{truth}}$, evaluated on the $Z\rightarrow\mu^{+}\mu^{-}$ simulated sample. Horizontal bars indicate the bin width; vertical bars show the 68\% Clopper-Pearson intervals, which are smaller than the markers. The drop at $\eta \approx 0$ for the Medium and Tight working points reflects the gap in the Muon Spectrometer coverage, while the Tight efficiency vanishes for $|\eta| > 2.5$, outside the Inner Detector acceptance.}
    \label{fig:efficiency}
\end{figure}
\begin{figure}[htbp]
  \centering
  \includegraphics[width=\textwidth]{TESI/purity_minimal.pdf}
  \caption{Purity of muons reconstructed as a function of (a) $p_\mathrm{T}$, (b) $\eta$ and (c) $\varphi$ for the Loose (blue), Medium (red) and Tight (green) identification working points, obtained from the simulated sample $Z\to\mu^+\mu^-$} 
  \label{fig:purezza_wp}
\end{figure}
\noindent A closer look in \cref{fig:muon_types} at the muon types explains these efficiencies. The most inclusive WP is the Loose, designed to maximize efficiency even in regions with acceptance holes. In fact, it admits all types of muons, but each only in the region where it is used to recover efficiency. In the range $|\eta| < 2.5$, CB muons dominate, accompanied by a lower IO contribution. The number of CBs drops sharply to $\eta \approx 0$, where the MS has a reduced coverage to make room for detector services, while in this region the IO contribution increases. For $|\eta| < 0.1$, the Loose also includes CT and ST muons, which do not require a complete trace in the spectrometer (CTs use the energy storage in the calorimeter, STs use a single segment in the MS) and therefore allow the identification of muons even where the MS has no coverage. Similarly, for $|\eta| > 2.5$, where the ID is no longer accepted, the Loose also selects SA and SiAF muons, extending the coverage up to $|\eta| = 2.7$. The WP Medium selects only the muons with momentum measured in the MS through a sufficient number of stations. For this reason, the CT and ST muons are filtered, which in the MS do not have a moment measurement (the CTs leave only an energy deposit in the calorimeter, while the STs leave a single segment in the MS). In the $|\eta| < 2.5$ muons are selected, which are required to hit at least two precision stations and the compatibility between the moment measured in the ID and the one measured in the MS. For $|\eta| < 0.1$, medium requires only one station, in order not to lose the muons that cross the gap due to the passage of services. Finally, for $|\eta| > 2.5$, where the ID is no longer accepted, the Medium accepts SA and SiAF muons, provided that they have hits in at least three precision stations, necessary to measure the curvature of the track using only the MS. The WP tight is defined within the acceptance of the ID, i.e. for $|eta|<2.5$, where it selects only the CB and IO muons, i.e. those with one track in the ID and one in the MS. To these it requires hits in at least two precision stations. After the Tight WP only CB and IO muons remain, within $|\eta|<2.5$, and the IO contribution is about two orders of magnitude smaller than the CB one. This explains the efficiency drop at $\eta\approx 0$ ($71\%$ for Medium, $60\%$ for Tight) and its vanishing for $|\eta|>2.5$. In \cref{tab:muon_types} are shown how many muons survive these working points criteria.

\begin{figure}[htbp]
    \centering
    \includegraphics[width=\textwidth]{TESI/muon_types_eta2.pdf}
    \caption{Pseudorapidity distribution of reconstructed prompt muons passing the (a) Loose, (b) Medium and (c) Tight working points in the $Z\rightarrow\mu^{+}\mu^{-}$ simulated sample, normalised to the number of events and stacked according to the reconstruction type: Combined, Inside-Out, Silicon-Associated Forward (SiAF), Stand-alone, Calo-Tagged and Segment-Tagged muons. Prompt muons are reconstructed muons matched to a truth muon classified as prompt. Calo-Tagged and Segment-Tagged muons enter only the Loose selection, in the region $|\eta| < 0.1$ not covered by the Muon Spectrometer, while Stand-alone and SiAF muons populate the region $2.5 < |\eta| < 2.7$ for the Loose and Medium working points. The Tight working point retains only Combined and Inside-Out muons within $|\eta| < 2.5$.}
    \label{fig:muon_types}
\end{figure}

\begin{table}[htbp]
  \centering
  \begin{tabular*}{\textwidth}{@{\extracolsep{\fill}}lrrr}
    \toprule
    Muon type & Loose & Medium & Tight \\
    \midrule
    Combined               & 325\,372\,621 & 325\,371\,313 & 312\,041\,086 \\
    Inside Out             & 2\,745\,809   & 2\,736\,673   & 1\,037\,356   \\
    Silicon-assoc. forward & 4\,754\,036   & 4\,754\,036   & 0             \\
    Stand-alone            & 4\,180\,914   & 4\,180\,914   & 0             \\
    Calo Tagged            & 3\,388\,486   & 0             & 0             \\
    Segment Tagged         & 73\,454       & 0             & 0             \\
    \bottomrule
  \end{tabular*}
  \caption{Number of muon types in the simulated $Z\to\mu^+\mu^-$ sample in the ideal geometry, that passing through Loose, Medium and Tight working points, associated to \cref{fig:muon_types}.}
  \label{tab:muon_types}
\end{table}

\clearpage

% =====================================================================

% =====================================================================
\section{Parametrization of the resolution}
\label{sec:parametrizzazione}

The trends obtained in the previous section, for each region of $\eta$, $p_{T}$ and $\varphi$, are the result of different mechanisms, each dominant in a different kinematic regime as previously discussed in Chapter 3. The parameterization of the resolution that we found in Eq. 3.14, we can now write it in the following way:
\begin{equation}
    \frac{\sigma(p_\mathrm{T})}{p_\mathrm{T}}= \frac{r_0}{p_\mathrm{T}} \oplus r_1 \oplus r_2\cdot p_\mathrm{T}
    \label{eq:parametrizzazione}
\end{equation}
The first term accounts for fluctuations in energy loss in the material being crossed and only matters at low momentum; the second, constant, accounts for multiple scattering and inhomogeneities in the description of the local magnetic field; the third, which
grows linearly with $p_{\mathrm{T}}$, describes spatial resolution
position measurements and the residual misalignments between the MS's chambers \cite{ATLAS-MCP-2023}. MS chamber misalignment produces two effect \cite{ATLAS-ID-ALIGN-2020}: a degradation of the track parameter resolution, which is even more noticeable the higher the momentum, and a parameter $\delta_{\mathrm{sagitta}}$ which alters the curvature in the opposite way for particles of different charge. The first one is the one that appears in $r_2$ and is the focus of this work. However, before going ahead with the analysis in the next sections, it is necessary to evaluate the behavior of the parameters reported in equation 7.2, through a correlation study with the coefficients $\rho$, using the following expression:
\begin{equation}
  \frac{\sigma(p_\mathrm{T})}{p_\mathrm{T}} = \sqrt{\left(\frac{r_0}{p_\mathrm{T}}\right)^2 + r_1^2 + \left(r_2\, p_\mathrm{T}\right)^2 + 2\rho_{01}\,\frac{r_0}{p_\mathrm{T}}\,r_1 + 2\rho_{02}\,r_0\, r_2 + 2\rho_{12}\, r_1\, r_2\, p_\mathrm{T}}
  \label{eq:fit_correlato}
\end{equation}
where $\rho_{01}$ indicates the correlations between $r_0$ and $r_1$, $\rho_{02}$ the correlation between $r_0$ and $r_2$ and $\rho{12}$ the correlation between $r_1$ ed $r_2$. By performing a quick fit to extrapolate the free parameters $\rho$, we get the following results reported in \cref{tab:rho} for all ranges of $|\eta^\mathrm{truth}|$. 
\begin{table}[htbp]
    \centering
    \begin{tabular*}{\textwidth}{@{\extracolsep{\fill}} c c c c c}
        \toprule
        $|\eta^{\mathrm{truth}}|$ &
        $\rho_{01}$ &
        $\rho_{02}$ &
        $\rho_{12}$ &
        $\chi^2/\mathrm{ndf}$ \\
        \midrule
        0.0--0.1  & $-0.3 \pm 0.5$ & $-0.7 \pm 1.7$ & $+1.0\pm 0.1$ & 0.18/10 \\
        0.1-1.01 & $-0.6 \pm 0.1$ & $-1$            & $+1$             & 0.16/10 \\
        1.01-1.3 & $+0.1 \pm 0.4$ & $-1$            & $+1$             & 0.67/10 \\
        1.3--1.7  & $-0.5 \pm 0.1$ & $-1$            & $+1$             & 0.92/10 \\
        1.7--2.5  & $-0.4 \pm 0.1$ & $-1$            & $+1$             & 0.41/10 \\
        2.5--2.8  & $+0.8 \pm 1.8$  & $-1.0 \pm 1.7$ & $-0.6 \pm 0.1$ & 12.9/9 \\
        \bottomrule
    \end{tabular*}
    \caption{Coerrelation coefficients $\rho_{01}$, $\rho_{02}$ and $\rho_{12}$ obtained from fitting $\sigma_{68}$ with eq.~7.3 in each interval of $|\eta^{\mathrm{truth}}|$ for the sample $Z\to\mu^+\mu^-$ in ideal geometry configuration. The values without uncertainty coincide with the limits of the allowed interval, $[-1, 1]$.}
    \label{tab:rho}
\end{table}
From this, we can see that the best way to parameterize equation 7.2 is with a linear expression of the three terms, according to the relationship:
\begin{equation}
    \frac{\sigma(p_\mathrm{T})}{p_\mathrm{T}}= \frac{r_0}{p_\mathrm{T}} + r_1 + r_2\cdot p_\mathrm{T}
\end{equation}
because with all three $\rho=\pm1$ we get a linear correlated relationship. As shown in \cref{fig:correlation}, where we took into account the most populated range $0.1<|\eta^\mathrm{truth}|<1.01$ for the sample $Z\to\mu^+\mu^-$ in an ideal geometry configuration, the linear fit trend (red) is the one that clearly follows the trend of equation 7.3 with the parameters left free (green). In blue is shown the fit performed through the relation:
\begin{equation}
    \frac{\sigma(p_\mathrm{T})}{p_\mathrm{T}} = \sqrt{\left(\frac{r_0}{p_\mathrm{T}}\right)^2 + r_1^2 + (r_2\ \cdot\ p_{\mathrm{T}})^2}
\end{equation}
a trend that would be obtained if the $\rho$ had been zero, and which clearly fails to adapt to the trend of the data. The graph shows the $\sigma_{68}$ interquantile that will be discussed in the next section.

\begin{figure}[htpb]
    \centering
    \includegraphics[width=\textwidth]{TESI/correlazione_minimal.pdf}
    \caption{$\sigma_{68}$ as a function of $p_\mathrm{T}^{\mathrm{truth}}$ in the range $0.10 < |\eta^{\mathrm{truth}}| < 1.01$ for the $Z\to\mu^+\mu^-$ sample in ideal geometry configuration, with the quadrature fit (blue), linear fit (red), and the fit of eq. 7.3 with free $\rho$ coefficients (green).}
    \label{fig:correlazione}
\end{figure}


\section{Alignment performance studies}

The $p_{\mathrm{T}}$ resolution analysis described by Eq. 7.1 is now repeated to compare two different configurations of the Muon Spectrometer: one in which all chambers are aligned (\emph{Ideal geometry}) and one in which the chambers are misaligned (\emph{MS misaligned geometry}, discussed in Section 7.3). For this purpose, two simulated samples, $Z$ and $Z'$, produced in both configurations, are used. Only prompt muons with momentum up to $\sim2$~TeV are considered, with the additional selection $\texttt{muon\_Type}==0$, which corresponds to CB muons only. As described in Chapter 6, Combined muons are indeed the only ones with a track in both the ID and the MS and, therefore, allow the effect of the chamber misalignment to be observed. Dividing the sample into 16 bins of $p_{\mathrm{T}}$ for each of the 6 bins of $|\eta|$, we can extract for each histogram the $68\%$ interquantile range defined in Eq.~7.5:
%
\begin{equation}
  \sigma_{68} = \frac{q_{84} - q_{16}}{2}\,,
  \label{eq:interquantile}
\end{equation}
%
This allow us to consider only the central area of the distribution, using the 16th and 84th percentiles, ignoring what happens in the tails. The error, instead, is just the statistical error $\sigma_{68}/\sqrt{2N}$. For each $|\eta|$ bin we can fit equation 
7.4 considering $\sigma_{68}(p_{\mathrm{T}}^{\mathrm{truth}})/p_{\mathrm{T}}$ performed on a log scale in \cref{fig:resq68-1}. The result of each parameters are shown in \cref{tab:risultati_fit_all}. Before discussing the results, a clarification is in order on why we chose to estimate the resolution through interquantiles. The distribution of the residual $\mathcal{R}$ is approximately symmetric around zero, since we expect the displacement introduced in the MS chambers to be uniform, so one might think of describing it with a Gaussian fit. However, as shown in \cref{fig:int_vs_gauss}, the distribution has a narrower peak and much longer tails than a Gaussian, resulting in a $\chi^2/\mathrm{ndf}$ much larger than one. The interquantile width $\sigma_{68}$, on the contrary, does not assume any shape for the distribution, since it only considers the central $68\%$ of the muon distribution, and it remains stable even in the presence of long tails. For this reason, we used it in the fits that follow.
\begin{figure}[htpb]
    \centering
    \includegraphics[width=\textwidth]{TESI/residuo_gauss_vs_q68_minimal.pdf}
    \caption{Distribution of the residual $\mathcal{R}$ for combined muons in the aligned simulated sample with $0.10 \le |\eta^{\mathrm{truth}}| < 1.01$, for $40\ \mathrm{GeV} \le p_{\mathrm{T}}^{\mathrm{truth}} < 42\ \mathrm{GeV}$ (a, b) and $500\ \mathrm{GeV} \le p_{\mathrm{T}}^{\mathrm{truth}} < 1000\ \mathrm{GeV}$ (c, d), in linear (a,c) and logarithmic (b, d) scale. The red curve is a Gaussian fit over the full range; the dashed lines mark the quantiles $q_{16}$ and $q_{84}$, and the shaded band the interval defining $\sigma_{68}$.}
    \label{fig:int_vs_gauss}
\end{figure}
\begin{figure}[htpb]
  \centering
  \includegraphics[width=\textwidth]{TESI/Fit_all_cb_minimal.pdf}
    \caption{Relative momentum resolution $\sigma_{68}/p_{\mathrm{T}}$ as a function of $p_{\mathrm{T}}^{\mathrm{truth}}$ for prompt muons of the combined $Z \to \mu^+\mu^-$ and $Z' \to \mu^+\mu^-$ samples in the ideal (perfectly aligned) geometry. Plots (a), (b), (c), (d), (e) and (f) show separately the six $|\eta|$ bins that follow the detector geometry. For each bin, $\sigma_{68}$ is extracted as half the difference between the 84th and 16th percentiles of the curvature distribution of \cref{eq:osservabile}, and the solid lines are the fits of \cref{eq:parametrizzazione}; the resulting parameters are listed in \cref{tab:risultati_fit_all}.}
  \label{fig:resq68-1}
\end{figure}
\noindent Returning, now to the alignment performance studies, the parameter $r_0$ is compatible with zero at low momentum, where the $r_0/p_{\mathrm{T}}$ term carries the most weight, because the CB track is dominated by the ID information. Since there is no material in front of the ID, energy-loss fluctuations are negligible and $r_0$ is therefore consistent with zero. In the region $2.5<|\eta|<2.8$ the ID no longer covers the acceptance (only a few residual CB muons remain), so the reconstruction is governed entirely by the MS. In this case, however, the MS resolution is affected by energy-loss fluctuations in the calorimeter located in front of it, which explains why the value in this region deviates significantly from zero ($0.108\pm0.018$~GeV) and leads to a large $\chi^2/\mathrm{ndf}$ in \cref{fig:resq68-1}~(f). The constant term $r_1$ increases with pseudorapidity, going from a minimum of $0.0139 \pm 0.0022$ in the central region of the detector to a maximum of $0.0256\pm0,0020$ in the transition region of the magnetic field. This trend matches the increase in material crossed and the local magnetic field irregularities that this term accounts for. The observed trend of $r_2$ is consistent with the bending power of the detector's magnet system, as shown in \cref{fig:bending_power} ,where $r_2\propto1/(BL^2)$. Indeed, the endcap ($1.7<|\eta|<2.5$), where the bending power is largest, yields a larger sagitta and consequently a small value of $r_2$, $(0.54\pm0.13)\cdot10^{-4}~\text{GeV}^{-1}$. Instead, the transition region, where the bending power is lowest, yields a smaller sagitta and a higher value of $r_2$, $(1.15\pm0.27)\cdot10^{-4}~\text{GeV}^{-1}$. The highest value of $r_2$, however, is observed in the central region ($|\eta|<0.1$), where it reaches $(1.33\pm0.31)\cdot10^{-4}~\text{GeV}^{-1}$. In this case the increase is not related to the bending power but to a geometric effect: a gap in the MS reserved for the passage of services~\cite{ATLAS-MUON-EFF-2021} degrades the measurement of the sagitta, resulting in a high value of $r_2$. The same region also shows the minimum value of the multiple scattering term $r_1$, since the muon traverses the least amount of material. The same effect was already observed in Section 7.3 as a drop in the efficiency of the Medium and Tight criteria.
\begin{figure}[htpb]
    \centering
    \includegraphics[width=\textwidth]{TESI/Toroid_field_integral.pdf}
    \caption{Bending power of the ATLAS toroidal magnet system, expressed as the field integral $\int B\,\mathrm{d}l$, as a function of $|\eta|$. The integral is evaluated along the trajectory of an infinite-momentum muon between the inner and outer MDT stations, for the azimuthal angles $\phi=0$ (red) and $\phi=\pi/8$ (black). The vertical lines delimit the transition region between barrel and endcap ($1.3<|\eta|<1.7$), where the bending power is lowest \cite{ATLAS-detector}.}
    \label{fig:bending_power}
\end{figure}

\begin{table}[htbp]
    \centering
    \begin{tabular}{c c c c c c}
        \toprule
        Plot & $|\eta^{\mathrm{truth}}|$ & $r_0$ [GeV] & $r_1$ & $r_2$ [$10^{-4}$ GeV$^{-1}$] & $\chi^2/\mathrm{ndf}$ \\
        \midrule
        (a) & $[0.00,\,0.10)$ & $-0.004 \pm 0.039$ & $0.0139 \pm 0.0022$ & $1.33 \pm 0.31$ & $1.79/13$ \\
        (b) & $[0.10,\,1.01)$ & $-0.040 \pm 0.031$ & $0.0174 \pm 0.0017$ & $0.90 \pm 0.22$ & $2.09/13$ \\
        (c) & $[1.01,\,1.30)$ & $\ \ 0.047 \pm 0.008$ & $0.0173 \pm 0.0005$ & $0.89 \pm 0.09$ & $1.08/13$ \\
        (d) & $[1.30,\,1.70)$ & $-0.019 \pm 0.034$ & $0.0256 \pm 0.0020$ & $1.15 \pm 0.27$ & $2.34/13$ \\
        (e) & $[1.70,\,2.50)$ & $-0.024 \pm 0.018$ & $0.0252 \pm 0.0010$ & $0.54 \pm 0.13$ & $1.98/13$ \\
        (f) & $[2.50,\,2.80)$ & $\ \ 0.108 \pm 0.018$ & $0.0205 \pm 0.0009$ & $0.62 \pm 0.12$ & $20.6/13$ \\
        \bottomrule
    \end{tabular}
    \caption{Parameters of the resolution fit of \cref{eq:parametrizzazione} in each $|\eta|$ bin, for the combined $Z$ and $Z'$ samples in the ideal (perfectly aligned) geometry, corresponding to the curves of \cref{fig:resq68-1}. The $\chi^2$ is evaluated with respect to the 13 degrees of freedom: 16 points in $p_{\mathrm{T}}$ minus the 3 free parameters.}
    \label{tab:risultati_fit_all}
\end{table}
\noindent As mentioned in Sec. 7.3, the choice of working point doesn't have a significant impact on the resolution. \cref{fig:resolution_wp} shows $\sigma_{68}$ as a function of $p_\mathrm{T}^\mathrm{truth}$ in different $|\eta^\mathrm{truth}|$ regions for the Loose, Medium, and Tight selections, for all types of muons. In the range $0.1 < |\eta^\mathrm{truth}| < 2.5$, the three curves overlap where both the ID and MS provide a measurement. So, the resolution is determined by the detector's performance and doesn't depend on the applied identification criteria. The effects of the selection only show up where one of the two detectors doesn't have coverage. The first of these regions is the central one ($|\eta^\mathrm{truth}| < 0.1$). Here, the MS has limited coverage, and the working points select muons with different measurement qualities to cover the missing acceptance. In particular, we’ve seen that the Loose WP also accepts muons whose momentum is basically determined by the ID, like CT and ST. The progressively stricter requirements of Medium and Tight, on the other hand, keep only the best-measured muons, leaving only CB and IO, but with a loss of efficiency as already noted in Sec. 7.2. The second one is the forward region shown in panel (f), $2.5 < |\eta^\mathrm{truth}| < 2.8$, which is beyond the ID acceptance and is only covered by the MS. In this region, there's also an unusual effect. As shown in \cref{fig:muon_types}, the Tight selection requires $|\eta^\mathrm{reco}| < 2.5$ on the reconstructed quantity. With perfect reconstruction, we wouldn’t expect Tight muons for $|\eta^\mathrm{truth}| > 2.5$, but here they appear to be present. To understand this effect, we need to take a step back and recall the role of the truth and reco quantities. In Monte Carlo simulation, the truth quantity is the one that is generated and is considered the absolute truth, meaning what we want to measure. The reco quantity, on the other hand, is what we actually measure after passing through the detector, and we don't know beforehand how accurate its reconstruction is. The residual compares the two quantities and allows us to quantify how well we are measuring. A residual close to 0 indicates that the reconstructed value is close to the truth one, as will be seen in the following sections. For a population of muons with a given momentum, the ideal would be a distribution of residuals as close as possible to a Dirac delta centered at 0. A shift of the center from 0 indicates a bias in the measurement, while the width of the distribution depends on the resolution with which the quantity is measured. Since neither the detector nor the reconstruction is perfect, any reconstructed quantity is affected by a resolution, and the simulation also takes this into account. This applies not only to $p_\mathrm{T}$, but also to $\eta$ and $\phi$. Also, as shown in \cref{fig:int_vs_gauss}, the resolution doesn’t follow a purely Gaussian pattern, but it has wide tails, which correspond to cases where the muon reconstruction is way off. The Tight muons observed for $|\eta^\mathrm{truth}| > 2.5$ are therefore muons for which the $\eta^\mathrm{reco}$ measurement was wrong, meaning their reconstructed $\eta$ is very different from the true one. Since the working point selections are based only on reconstructed quantities, these muons therefore have $|\eta^\mathrm{reco}| < 2.5$ and pass the Tight selection, even though they have $|\eta^\mathrm{truth}| > 2.5$. Reconstruction errors this big can originate in the detector, for example due to local efficiency losses that make you miss hits crucial for the reconstruction. But they can also depend on the software, since fitting the muon track is particularly tricky. Figuring out which of the two is responsible isn’t easy and goes beyond the scope of this work.
\begin{figure}[htpb]
    \centering
    \includegraphics[width=0.85\textwidth]{TESI/Fit_all_wp_etatruth_minimal.pdf}
    \caption{Curvature resolution $sigma_{68}$ as a function of $p_\mathrm{T}^\mathrm{truth}$ for i working points Loose (blue), Medium (red) and Tight (green), in the regions of $|\eta^\mathrm{truth}|$ indicated in panels (a)-(f).}
    \label{fig:resolution_wp}
\end{figure}


\clearpage

\section{Residual Misalignment performance studies}

With the same identical Section 7.5 procedures, we can now analyze the MS misaligned geometry samples of $Z$ and $Z'$, as shown in \cref{fig:q68MS} and the result of each parameters are shown in \cref{tab:risultati_fit_disall}. The values of $r_0$ and $r_1$ do not depend on the alignment of the MS chambers and show no significant variations between the two configurations. What is instead sensitive to a different configuration of the MS chambers is the $r_2$ term, since it depends on the precision of the sagitta measurement, and it is systematically larger than in the ideal configuration, reflecting the introduced misalignment. 
\begin{figure}[htpb]
  \centering
  \includegraphics[width=\textwidth]{TESI/Fit_disall_cb_minimal.pdf}
    \caption{Relative momentum resolution $\sigma_{68}/p_{\mathrm{T}}$ as a function of $p_{\mathrm{T}}^{\mathrm{truth}}$ for prompt muons of the combined $Z \to \mu^+\mu^-$ and $Z' \to \mu^+\mu^-$ samples in the MSmisalign geometry, where the chambers of the middle layer of the Muon Spectrometer are displaced by amounts compatible with the residual precision of the alignment procedure, shown separately for the six $|\eta|$ bins. As in \cref{fig:resq68-1}, $\sigma_{68}$ is extracted from the percentiles of the curvature distribution and the solid lines are the fits of \cref{eq:parametrizzazione}; the resulting parameters are listed in \cref{tab:risultati_fit_disall}. The overall resolution is degraded with respect to the ideal geometry and the degradation grows with $p_{\mathrm{T}}$, as expected for a term proportional to $p_{\mathrm{T}}$.}
  \label{fig:q68MS}
\end{figure}
\begin{table}[htbp]
    \centering
    \begin{tabular*}{\textwidth}{@{\extracolsep{\fill}} c c c c c c @{}}
        \toprule
        Plot & $|\eta^{\mathrm{truth}}|$ & $r_0$ [GeV] & $r_1$ & $r_2$ [$10^{-4}$ GeV$^{-1}$] & $\chi^2/\mathrm{ndf}$ \\
        \midrule
        (a) & $[0.00,\,0.10)$ & $0.017 \pm 0.040$ & $0.0127 \pm 0.0024$ & $2.20 \pm 0.35$ & $0.53/13$ \\
        (b) & $[0.10,\,1.01)$ & $0.014 \pm 0.026$ & $0.0153 \pm 0.0016$ & $1.54 \pm 0.23$ & $0.89/13$ \\
        (c) & $[1.01,\,1.30)$ & $0.069 \pm 0.023$ & $0.0160 \pm 0.0015$ & $1.90 \pm 0.23$ & $1.16/13$ \\
        (d) & $[1.30,\,1.70)$ & $0.056 \pm 0.023$ & $0.0214 \pm 0.0015$ & $2.05 \pm 0.23$ & $2.75/13$ \\
        (e) & $[1.70,\,2.50)$ & $0.070 \pm 0.016$ & $0.0203 \pm 0.0011$ & $1.82 \pm 0.18$ & $2.94/13$ \\
        (f) & $[2.50,\,2.80)$ & $0.217 \pm 0.046$ & $0.0148 \pm 0.0026$ & $2.54 \pm 0.37$ & $9.98/12$ \\
        \bottomrule
    \end{tabular*}
    \caption{Parameters of the resolution fit of \cref{eq:parametrizzazione} in each $|\eta|$ bin, for the combined $Z$ and $Z'$ samples in the MS misaligned geometry, corresponding to the curves of \cref{fig:q68MS}. The $\chi^2$ is evaluated with respect to the 13 degrees of freedom: 16 points in $p_{\mathrm{T}}$ minus the 3 free parameters.}
    \label{tab:risultati_fit_disall}
\end{table}
\noindent The residual misalignment term can therefore be extracted by subtracting in quadrature the values of $r_2$ obtained in the configuration with all chambers aligned from those obtained in the configuration with misaligned MS chambers, as shown in Eq. 6.6:
\begin{equation}
  r_2^{\mathrm{residual}} = \sqrt{\left(r_2^{\mathrm{MSmisalign}}\right)^2 -
  \left(r_2^{\mathrm{Aligned}}\right)^2}
  \label{eq:residui-disallineamento}
\end{equation}
The uncertainties on $r_2^{\mathrm{residual}}$ are evaluated by using the central $68\%$ interval of the distribution. This allows us to understand how much alignment impacts the resolution. The results are shown in \cref{tab:result_residual}.
\begin{table}[htbp]
  \centering
  \begin{tabular*}{\textwidth}{@{\extracolsep{\fill}} clccc @{}}
    \toprule
    & $|\eta^{\mathrm{truth}}|$
    & $r_2^{\mathrm{Aligned}}$ $[10^{-4}\,\mathrm{GeV}^{-1}]$
    & $r_2^{\mathrm{MSmisalign}}$ $[10^{-4}\,\mathrm{GeV}^{-1}]$
    & $r_2^{\mathrm{residual}}$ $[10^{-4}\,\mathrm{GeV}^{-1}]$ \\
    \midrule
    (a) & $[0.00,\,0.10)$ & $1.33 \pm 0.31$ & $2.20 \pm 0.35$ & $1.74^{+0.45}_{-0.59}$ \\
    (b) & $[0.10,\,1.01)$ & $0.90 \pm 0.22$ & $1.54 \pm 0.23$ & $1.26^{+0.29}_{-0.38}$ \\
    (c) & $[1.01,\,1.30)$ & $0.89 \pm 0.09$ & $1.90 \pm 0.23$ & $1.67^{+0.26}_{-0.27}$ \\
    (d) & $[1.30,\,1.70)$ & $1.15 \pm 0.27$ & $2.05 \pm 0.23$ & $1.70^{+0.30}_{-0.38}$ \\
    (e) & $[1.70,\,2.50)$ & $0.54 \pm 0.13$ & $1.82 \pm 0.18$ & $1.74^{+0.19}_{-0.20}$ \\
    (f) & $[2.50,\,2.80)$ & $0.62 \pm 0.12$ & $2.54 \pm 0.37$ & $2.47^{+0.38}_{-0.39}$ \\
    \bottomrule
  \end{tabular*}
    \caption{Values of $r_2$ obtained from the fit with all chambers aligned
    ($r_2^{\mathrm{Aligned}}$) and with misaligned MS chambers
    ($r_2^{\mathrm{MSmisalign}}$), and residual misalignment term
    $r_2^{\mathrm{residual}}$ extracted with Eq.~6.6,
    for each $|\eta^{\mathrm{truth}}|$ region. Only combined muons are used.
    The uncertainties on $r_2^{\mathrm{residual}}$ correspond to the $68\%$
    interval of the distribution.}
    \label{tab:result_residual}
\end{table}
\noindent Another way to visualize how much the resolution changes with and without misalignment is shown in \cref{fig:res_cb}. When all muons are considered ($20\ \mathrm{GeV} \le p_{\mathrm{T}} < 2000\ \mathrm{GeV}$) the effect is small, whereas at high transverse momentum ($p_{\mathrm{T}} \ge 200$~GeV) the uniform misalignment of the MS chambers significantly degrades the resolution.
\begin{figure}[htpb]
    \centering
    \includegraphics[width=0.75\textwidth]{TESI/residuo_allineato_vs_disallineato_minimal.pdf}
    \caption{Distribution of the residual $\mathcal{R}$ for combined muons with
  $0.10 \le |\eta^{\mathrm{truth}}| < 1.01$ in the aligned sample (blue) and in the sample
  with misaligned MS chambers (red), for $20\ \mathrm{GeV} \le p_{\mathrm{T}}^{\mathrm{truth}} < 2000\ \mathrm{GeV}$ (a) and $200\ \mathrm{GeV} \le p_{\mathrm{T}}^{\mathrm{truth}} < 2000\ \mathrm{GeV}$ (b). The corresponding values of $\sigma_{68}$ are given in the legend.}
    \label{fig:res_cb}
\end{figure}
\noindent To summarize, the two configurations use the same samples and the same selection, differing only in the alignment of the MS chambers. Any systematic difference can therefore be attributed solely to this misalignment. The $r_2$ term increases in all 6 bins in the MS misaligned configuration with respect to the ideal one. From \cref{tab:result_residual} it can be observed that the residual term is practically constant in bins (a) to (e), compatible within the uncertainties. This is exactly what is expected if the misalignment introduces a uniform contribution to the sagitta measurement, independent of the region. In bin (f) the contribution is higher because, in the absence of the ID, the sagitta is measured only by the MS and only a few CB muons can be adequately reconstructed.


% ==================================================================
%  CAPITOLO 6 - CONCLUSIONI
% ==================================================================
\chapter{Conclusion}
% Testo originale conservato come traccia (era gia' commentato nel file di partenza):
%Muon reconstruction and identification in the ATLAS detector are critical processes to ensure the accuracy of physics analyses, particularly in rejecting background from non-prompt sources. [...]

\begin{todo}[Conclusioni da scrivere]
Da fare: riassumere i risultati ottenuti -- validazione del matching truth/reco,
frazione di muoni prompt (99.25\%), andamento della risoluzione in \(p_T\) e in
\(|\eta|\), parametri \(r_0\), \(r_1\), \(r_2\) e confronto con i campioni disallineati --
e indicare i possibili sviluppi futuri.
\end{todo}


\bibliographystyle{unsrtnat}
\bibliography{references}     % senza estensione .bib
\end{document}